\documentclass[11pt]{article}
\usepackage[T1]{fontenc}
\usepackage{lmodern}
\usepackage[margin=1in]{geometry}
\usepackage{amsmath,amssymb,amsthm,mathtools,amscd}
\usepackage{enumitem}
\usepackage[hidelinks]{hyperref}
\usepackage{microtype}
\hypersetup{pdftitle={Topological triviality of mu-constant families of surface singularities},pdfauthor={OpenAI}}
\pdfinfoomitdate=1
\pdftrailerid{}
\pdfsuppressptexinfo=15
\newtheorem{theorem}{Theorem}[section]
\newtheorem{proposition}[theorem]{Proposition}
\newtheorem{lemma}[theorem]{Lemma}
\newtheorem{corollary}[theorem]{Corollary}
\theoremstyle{definition}
\newtheorem{definition}[theorem]{Definition}
\theoremstyle{remark}
\newtheorem{remark}[theorem]{Remark}
\DeclareMathOperator{\Supp}{Supp}
\DeclareMathOperator{\Sing}{Sing}
\DeclareMathOperator{\Spec}{Spec}
\DeclareMathOperator{\edim}{edim}
\title{Topological triviality of $\mu$-constant families of surface singularities}
\author{OpenAI}
\date{September 24, 2026}
\begin{document}
\maketitle
\begin{abstract}
We prove that every holomorphic one-parameter family of isolated hypersurface
singularities in $\mathbb C^3$ with constant Milnor number is topologically
right-trivial. This gives a positive answer to the surface case of the
$\mu$-constant problem. The trivialization fixes the parameter and the
origin section, and preserves the defining functions.
\end{abstract}
\tableofcontents
\section{Introduction}\label{sec:introduction}

The Milnor number measures the topology of a nearby smooth level of an
isolated hypersurface singularity: its Milnor fiber has the homotopy type
of a bouquet of spheres, with one sphere for each unit of the Milnor
number \cite[Theorems~6.5 and~7.2]{Milnor1968}; the analytic
formulation in terms of the Jacobian algebra is recalled in
\cite[Section~1, Theorem~1.5]{LeRamanujam1976}.  A basic question in
equisingularity theory asks whether keeping this number constant also
keeps the defining function topologically unchanged.  For complex
surfaces, this requires a method that works beyond the dimensions in
which the classical cobordism argument applies.

Let $\Delta\subset\mathbb C$ be a disk centered at zero and let
\[
 F:(\mathbb C^3\times\Delta,\{0\}\times\Delta)\longrightarrow(\mathbb C,0)
\]
be holomorphic.  Write $f_t(x)=F(x,t)$ and suppose that $f_t$ has an
isolated critical point at the origin.  Its Milnor number is
\[
 \mu(f_t,0)=\dim_{\mathbb C}
 \frac{\mathcal O_{\mathbb C^3,0}}
 {(\partial f_t/\partial x_1,\partial f_t/\partial x_2,
   \partial f_t/\partial x_3)}.
\]
The family is \emph{topologically right-trivial} if a homeomorphism germ
of the ambient family, preserving the parameter and the origin section,
identifies each defining function with the central one.  The coordinate
changes and their inverses must vary jointly continuously with the
parameter.

\begin{theorem}\label{thm:main}
Let $F:(\mathbb C^3\times\Delta,\{0\}\times\Delta)\to(\mathbb C,0)$
be holomorphic.  Suppose that, for every $t\in\Delta$, $f_t(0)=0$,
$f_t$ has an isolated critical point at $0$, and
$\mu(f_t,0)=\mu(f_0,0)$.  After shrinking the disk and the germ
representatives, there is a homeomorphism germ
\[
 \Phi(x,t)=(\phi_t(x),t),\qquad
 \phi_t(0)=0,\quad \phi_0=\mathrm{id},\qquad
 f_t\bigl(\phi_t(x)\bigr)=f_0(x).
\]
Both $\Phi$ and $\Phi^{-1}$ are jointly continuous on neighborhoods
of the origin section over the original disk.
\end{theorem}

Thus the surface case of the $\mu$-constant problem has a positive
answer.  The parameter dependence may be arbitrary and holomorphic, and
the surface singularities need not be log canonical.  The finite
parameter covers used in the proof are removed before constructing the
homeomorphism in the theorem.  The same ambient maps identify the zero
sets and the nearby nonzero levels, keeping the value of the function
fixed throughout the parameter disk.

\subsection*{The problem and its antecedents}

L\^e and Ramanujam proved that a $\mu$-constant family of isolated
hypersurface singularities has constant embedded topological type when
the hypersurface dimension is different from two
\cite[Theorem~2.1]{LeRamanujam1976}.  Their argument compares Milnor
fibers through a homologically trivial cobordism.  In higher dimensions
the $h$-cobordism theorem turns this comparison into a product; the
curve case follows from surface topology.  Complex dimension two is
the dimension in which this argument leaves the product structure
unresolved.  Outside complex dimension two, Timourian and King established local
triviality of families of function germs
\cite{Timourian1977,King1980}.  This additional passage matters here:
constancy of the individual topological types and a jointly continuous
choice of ambient coordinate changes are different conclusions.

Whitney regularity supplies a stronger geometric form of equisingularity.
For a hypersurface germ in three variables, let $\mu^{(j)}$ be the
Milnor number of its restriction to a general $j$-dimensional linear
subspace.  Then $\mu^{(3)}=\mu$, while $\mu^{(1)}$ is the multiplicity
minus one and $\mu^{(2)}$ records the generic plane section.  Teissier
proved that constancy of this full sequence implies Whitney regularity
along the singular section \cite[Chapter~II, Theorem~3.9]{Teissier1973};
Brian\c con and Speder proved the converse
\cite[Theorem~2]{BrianconSpeder1976}.  Their examples also show that
topological triviality need not imply Whitney regularity
\cite{BrianconSpeder1975}\cite[Section~3]{BrianconSpeder1976}.
Thus a proof under the sole $\mu$-constant hypothesis must account for
families whose generic plane-section Milnor number changes.

The equimultiplicity theorem of Fern\'andez de Bobadilla and Pe\l ka
settles a separate part of this numerical comparison: a family with
constant finite Milnor number has constant multiplicity in every
dimension \cite[Theorem~1.1]{FernandezBobadillaPelka2024}.  Their theorem
allows even formal power series whose coefficients vary continuously
along a real interval.  Its proof gives A'Campo's radius-zero
monodromy model a symplectic structure and extends McLean's
Floer-theoretic spectral sequence to compare multiplicities without
assuming a common Milnor radius.  For the present surface problem it supplies constancy of
$\mu^{(1)}$, but does not supply constancy of $\mu^{(2)}$.  We use it
with Teissier's criterion in the log-canonical case; the double-point
case reduces to the classical $\mu$-constant problem for plane curves.

Vector fields offer another approach to the parameter dependence.
The L\^e--Saito--Teissier criterion relates constancy of the Milnor
number to the smallness of the parameter derivative compared with the
spatial gradient \cite{LeSaito1973}\cite[Chapter~II, Remark~3.10]{Teissier1973}.
Such an estimate makes a gradient lift continuous at the singular
section, but continuity alone does not give unique, continuously
invertible transport.  Parusi\'nski constructed trivializing vector
fields for $\mu$-constant linear families $f(x)+tg(x)$ of isolated
hypersurface singularities, without a dimension restriction
\cite[Corollary~2.1]{Parusinski1999}.  Appendix~\ref{sec:continuous-gradient}
separates the gradient estimate from the trajectory control needed
for general parameter dependence.

For surfaces, simultaneous resolution gives a bridge from numerical
invariants to topology.  Wahl associated a logarithmic characteristic
number to a normal surface singularity \cite{Wahl1990}.
Okuma showed that, for a deformation of a non-log-canonical Gorenstein
surface singularity, constancy of this number produces simultaneous surface resolutions
after finite base change, with reduced exceptional divisors on the
fibers over the distinguished singular section
\cite[Theorem~1.2]{Okuma2004}.

Recent work describes further constraints on degenerations under the
$\mu$-constant hypothesis.  Aldasoro Rosales studies semistable models
after suitable base change and projective compactification.  Her
results include rational acyclicity of the special-fiber dual complex,
Tate type for the second cohomology of the additional vertical
surfaces, and equality of the first Betti numbers of the special fiber
and its distinguished component
\cite[Theorem~1.2]{AldasoroRosales2025}.  The comparison below likewise
studies these additional surfaces, using their boundary complements
and logarithmic canonical squares.  Peters proves symplectic
equivalence of completed Milnor fibers and contact equivalence of
links for $\mu$-constant polynomial families without vanishing folds
\cite[Theorem~3.7]{Peters2026}.  That extra hypothesis provides a common
Milnor radius.  Our argument instead obtains transport from a resolved
zero family and extends it to the ambient function levels.

\subsection*{The resolution invariant and the numerical comparison}

Write $X_t=\{f_t=0\}$.  These are normal Gorenstein surfaces.  For
non-log-canonical germs, our main numerical task is to deduce constancy
of Wahl's logarithmic surface invariant from $\mu$ alone.
On a good resolution, write the relative Zariski decomposition of the
canonical divisor plus the full reduced exceptional divisor as $P+N$,
with $P$ relatively nef and $N$ effective and exceptional.  The
exceptional intersection matrix identifies a unique exceptional rational
cycle $P_{\mathrm{loc}}$ having the same intersections as $P$ with every
exceptional curve.  The invariant is
\[
 \beta=-P_{\mathrm{loc}}^2.
\]
Section~\ref{sec:local-square} recalls its resolution independence and
the conventions used by Okuma \cite[Section~2]{Okuma2004}.  It also
expresses this local square as a difference of ordinary projective
squares.

The proof rests on a cancellation of Euler characteristics.  To explain
it, first consider a projective family with just the prescribed surface
singularity in each fiber.  A semistable log minimal model replaces its
special fiber by a component $S_0$ dominating $X_0$ and additional
projective surfaces $S_a$ over the singular point.  Each component $S$
has a reduced boundary $B_S$, consisting of its intersections with the
other vertical components and the horizontal exceptional divisors.
Write $S^\circ=S\setminus B_S$.  The construction gives
$S_0^\circ\simeq X_0\setminus\{0\}$ and makes every
$K_{S_a}+B_{S_a}$ nef.

Constancy of $\mu$ keeps the Euler characteristic of the punctured
projective fibers constant.  Their nearby Euler comparison, after the
contribution of $S_0^\circ$ cancels, is
\[
 0=\sum_a\chi_c(S_a^\circ)+\delta,\qquad \delta\ge0.
\]
Here $\chi_c$ denotes Euler characteristic with compact supports, and
$\delta$ is the sum of the excess Euler characteristics of the local
smoothings at interior singular points of the extra surfaces.
Greuel--Steenbrink's first-Betti-number theorem makes these excesses
nonnegative \cite{GreuelSteenbrink1983}.
Langer's orbifold Bogomolov--Miyaoka--Yau inequality supplies the other
sign \cite{Langer2003}:
\[
 0\le (K_{S_a}+B_{S_a})^2\le3\chi_c(S_a^\circ)
 \qquad\text{for every }a.
\]
Thus every extra logarithmic square vanishes.  Specialization of
projective intersections then gives constancy of $\beta$.  The key
point is to retain a boundary for which both the Euler comparison and
the surface inequality hold.

The technical development establishes these properties in order.
Section~\ref{sec:algebraic-reduction} reduces this scalar assertion to
polynomial families over algebraic curves.  Individual finite
determinacy supplies polynomial representatives, and constructibility
of the resolution invariant gives the curve comparison.  These
individual equivalences are used only to compare numerical invariants.
The high-degree compactification, following the construction in
\cite[Lemma~2.5]{Steenbrink1985}, gives projective surfaces with precisely
the prescribed singularity and a smooth complement.

Sections~\ref{sec:log-model}--\ref{sec:full-boundary} use semistable
reduction and the threefold log minimal model program
\cite[Chapter~7]{KM1998} to obtain a model with a reduced special
fiber.  The boundary retains every surviving horizontal exceptional
divisor.  Relative nefness and the negativity lemma identify its full
support over the singular section.  Adjunction then supplies precisely
the reduced boundaries used in the comparison above.

Sections~\ref{sec:nearby-euler}--\ref{sec:surface-estimates} establish
the Euler comparison and its signs.  The essential local assertion is
that a boundary intersection contributes zero to the nearby Euler
function, including at singular points of the threefold.  Finite covers
that make the boundary Cartier and the structure of threefold dlt strata reduce
this assertion to the Euler calculation for a product of coordinates.
Proper compatibility of nearby cycles then permits the local
calculation to be integrated over the special fiber
\cite{SchuermannNearby,MaximSchuermann}.  The interior correction and the
comparison of orbifold and ordinary Euler characteristics complete the
two sign calculations.  Section~\ref{sec:square-specialization} converts
the resulting vanishing into constancy of $\beta$.

The log-canonical branch is handled separately.  Constancy of the
spectrum, following Varchenko and Steenbrink
\cite{Varchenko1982,Steenbrink1985}, together with the threshold formula
\cite[Theorem~9.5]{Kollar1997Pairs} and inversion of adjunction
\cite{Kawakita2007}, shows that log canonicity is constant along the
family.  Section~\ref{sec:low-cases} treats this branch and the
double-point case through sectional Milnor numbers.

\subsection*{From resolution to a homeomorphism of the functions}

In the non-log-canonical branch, Okuma's theorem produces a simultaneous
minimal semigood resolution
after a finite disk cover \cite[Theorem~1.2]{Okuma2004}.  Its reduced
exceptional restrictions are essential: they allow the remaining
relative-normal-crossings assertion to be proved by the Euler
characteristic of the exceptional curves.  A node smoothing would
lower that characteristic, whereas the resolved tubes and the
punctured zero fibers make it constant.

Section~\ref{sec:ambient-lifting} constructs complete lifts of the
parameter directions on the resolved tubes.  These descend to the
punctured zero fibers and extend to the ambient nonzero levels while
preserving the defining function.  Properness of the tubes and reverse
transport establish joint continuity at the origin for both the maps
and their inverses.  An explicit correction along the angular seam of
$t=u^e$ then yields a trivialization over the original disk.
Section~\ref{sec:low-cases} supplies the double-point and log-canonical
cases using sectional Milnor numbers and completes the proof.


Appendix~\ref{app:local-numerical} supplies the marked monomial
construction used in Section~\ref{sec:log-model}.  It also gives two
variants of the numerical argument: specialization of intersection
squares supported properly over a curve and comparison within an
irreducible component of a finite jet space.
Appendix~\ref{sec:continuous-gradient} gives an ambient transport
criterion using continuous gradient lifts and explicit control of
trajectories at the singular section.  Its two sufficient geometric
conditions come from Whitney regularity and from a proper simultaneous
resolution with relative normal crossings.  The criterion separates the
smallness of the vector fields from the additional control needed to
obtain continuous flows and continuous inverse flows.

\section{Local families and fixed representatives}
\label{sec:local-families}

We first record the local information needed for the numerical comparison
and for the eventual ambient trivialization.  Throughout this section,
$F(x,t)=f_t(x)$ is a holomorphic family in three spatial variables over a
disk $\Delta$, with $f_t(0)=0$, an isolated critical point at $0$, and
$\mu(f_t,0)=\mu>0$ for every $t\in\Delta$.  Write
$X=F^{-1}(0)$, $X_t=f_t^{-1}(0)$, and
$\Gamma=\{0\}\times\Delta$.  Representatives and the parameter disk may
always be shrunk about the central point.

\begin{lemma}[Local family properties]
\label{lem:local-family-properties}
There are representatives on which the spatial critical locus of $F$ is
exactly $\Gamma$.  On these representatives $X\to\Delta$ is flat, its
fibers are normal Gorenstein surfaces, and $X$ is a normal Gorenstein
threefold, smooth over $\Delta$ away from $\Gamma$.  The same statements
hold after a nonconstant holomorphic base change between smooth disks.
\end{lemma}

\begin{proof}
Consider the scheme defined by the three spatial partial derivatives.
Its central fiber is supported at $0$ and has length $\mu$.  The local
finiteness theorem for analytic maps gives a finite representative over
the disk.  In the regular local ring $\mathbb C\{x_1,x_2,x_3,t\}$, the
three partials together with $t$ form a system of parameters.  They
therefore form a regular sequence.  The quotient by the three partials
is Cohen--Macaulay of dimension one, and $t$ is a nonzerodivisor in it.
It is consequently finite flat over the disk, of degree $\mu$.
In every nearby fiber its local length at $0$ is already $\mu$; there
can be no further points.  This is conservation of the spatial-gradient
number; compare \cite[Proposition~3.1]{Denkowski2018}.

Each surface fiber is a hypersurface, hence Gorenstein and
Cohen--Macaulay.  Its isolated singularity has codimension two, so
Serre's criterion gives normality.  Similarly, the total hypersurface
is Gorenstein and Cohen--Macaulay; it is smooth off $\Gamma$, which has
codimension two in $X$, and is therefore normal.  Since no fiber of the
ambient projection is contained in $X$, its parameter is a
nonzerodivisor on $X$.  Torsion-freeness over the regular one-dimensional
base gives flatness.  After a disk base change the equation is again a
relative hypersurface, with the same individual surface germs and the
same spatial critical locus.  These arguments apply again.
\end{proof}

Use $\chi_c$ for Euler characteristic with compact supports; on a compact
space it agrees with ordinary Euler characteristic $\chi$.

\begin{lemma}[Euler characteristic in a fixed ball]
\label{lem:fixed-ball-euler}
There are a sufficiently small closed Euclidean ball $B$ about $0$, a
smaller parameter disk, and $\eta>0$ such that the following statements
hold for every parameter in that disk:
\begin{enumerate}
\item $0$ is the only spatial critical point in $B$;
\item $f_t^{-1}(c)$ meets $\partial B$ transversely whenever $|c|<\eta$;
\item
\[
 \chi(X_t\cap B)=1,
 \qquad
 \chi_c\bigl((X_t\cap B)\setminus\{0\}\bigr)=0.
\]
\end{enumerate}
The Euler equalities also hold if the radius of $B$ varies slightly
within a sufficiently small interval about the chosen radius.
\end{lemma}

\begin{proof}
Choose $B$ inside the representative of
Lemma~\ref{lem:local-family-properties}, with radius a Milnor radius for
$f_0$.  Transversality at the compact set
$f_0^{-1}(0)\cap\partial B$ persists for small $t$ and small $c$.
Points elsewhere on the sphere do not enter these small levels, after
shrinking, by compactness.  This proves the first two assertions.

Fix $c$ with $0<|c|<\eta$.  The spaces $f_t^{-1}(c)\cap B$ form a proper
smooth family with boundary over the smaller parameter disk.  Smooth
lifts tangent to the boundary give the proper-submersion theorem in
this setting.  Their Euler characteristic is therefore that of the
central Milnor fiber, namely $1+\mu$; see
\cite[Theorems~6.5 and~7.2]{Milnor1968} and, for the analytic
Jacobian-algebra formulation, \cite[Section~1, Theorem~1.5]{LeRamanujam1976}.

For a fixed $t$, choose an inner Milnor ball for $f_t$, with closure in
the interior of $B$.  On the closed region between the two spheres,
$f_t$ has no critical point, and both boundary restrictions are
submersive over a sufficiently small disk of values.  Thus the shell
has the same Euler contribution at value $0$ and at sufficiently small
nonzero values.  In the inner ball, passing from value $0$ to a small
nonzero value replaces a contractible singular fiber by a Milnor fiber
of Euler characteristic $1+\mu$.  Additivity of $\chi_c$, or a
decomposition using the inner boundary and the two complementary
pieces, therefore gives
\[
 \chi(f_t^{-1}(c)\cap B)-\chi(X_t\cap B)=\mu.
\]
The inner ball and the admissible nonzero value may depend on $t$;
the preceding computation of the outer smooth level applies to every
such value.  It follows that $\chi(X_t\cap B)=1$.  Removing its one
origin point gives the second equality.  Finally, transversality
persists for radii in a small interval around the chosen radius, and
the identical argument applies there.
\end{proof}

Our algebraic reduction will concern the numerical invariant of each
individual surface germ.  The following elementary form of finite
determinacy supplies polynomial representatives for that purpose.

\begin{lemma}[Individual finite determinacy]
\label{lem:individual-jets}
Let $f\in\mathbb C\{x_1,x_2,x_3\}$ have an isolated critical point at $0$
and Milnor number $\mu$, and write $\mathfrak m=(x_1,x_2,x_3)$.
For every $b\in\mathfrak m^{\mu+2}$ there is a biholomorphic germ $h$
fixing $0$ such that $(f+b)\circ h=f$.  In particular, if
$N\geq\mu+1$, then $f$ is analytically right-equivalent to its Taylor
polynomial $j^Nf$.
\end{lemma}

\begin{proof}
Let $\mathfrak m=(x_1,x_2,x_3)$ and $J(f)=(f_{x_1},f_{x_2},f_{x_3})$.
The local algebra $\mathbb C\{x\}/J(f)$ has length $\mu$, so its maximal
ideal has $\mu$th power zero, by the descending filtration and
Nakayama's lemma.  Hence $\mathfrak m^\mu\subset J(f)$.
For the given $b\in\mathfrak m^{\mu+2}$ we have
\[
 b\in\mathfrak m^{\mu+2}\subset\mathfrak m^2J(f),
 \qquad
 b_{x_i}\in\mathfrak m^{\mu+1}\subset\mathfrak m J(f).
\]
Write $b=a^{\mathsf T}g$, where $g=(f_{x_i})_i$ is a column vector and
the entries of $a$ lie in $\mathfrak m^2$.  There is a matrix $C$ with
entries in $\mathfrak m$ such that the gradient column of $f+sb$ is
$(I+sC)g$.  On a sufficiently small neighborhood this matrix is
invertible simultaneously for $0\leq s\leq1$.  The holomorphic spatial
vector field
\[
 w_s=-(I+sC)^{-\mathsf T}a
\]
vanishes to second order and satisfies $d(f+sb)(w_s)=-b$.
Its time-dependent flow $h_s$, initially the identity, exists as a
holomorphic germ throughout $0\leq s\leq1$, after shrinking the spatial
neighborhood.  This follows from a uniform quadratic bound on $w_s$;
the reverse equation supplies the inverse germ.  The chain rule gives
$\frac{d}{ds}((f+sb)\circ h_s)=0$, and therefore $h=h_1$ has the
required property.  Taking $b=j^Nf-f$ proves the assertion about Taylor
polynomials.
\end{proof}

\subsection{Birational conventions}
A pair $(Z,B)$ consists of a normal complex variety and an effective rational
Weil divisor $B$ such that $K_Z+B$ is rational Cartier.  For a prime divisor
$E$ on a birational model, its log discrepancy is one plus the coefficient
of $E$ in the difference between the canonical divisor upstairs and the
pullback of $K_Z+B$, using compatible canonical divisors.  The pair is
\emph{log canonical}, abbreviated lc, if all log discrepancies are nonnegative.
For $B=0$ we say that $Z$ is log canonical.  A pair with boundary coefficients
at most one is \emph{divisorial log terminal}, abbreviated dlt, if it has a
log resolution on which every exceptional divisor has positive log
discrepancy.  An lc center is the image of a divisor with log discrepancy zero.
For an lc pair, the equivalent dlt criterion is that a simple
normal crossing open set contains the generic points of all lc centers.
We use the usual normality and adjunction properties of its coefficient-one strata
\cite[11.5--11.15]{Kollar2023Families}.
A normal variety is $\mathbb Q$-factorial if every prime Weil divisor has a
positive Cartier multiple.  A divisor is nef over a morphism when it has
nonnegative degree on every complete contracted curve.

\section{The local logarithmic square}
\label{sec:local-square}

The numerical invariant used below can be computed on a partial resolution.
We first specify its local intersection convention, then express it as a
difference of ordinary projective intersection numbers.

Let $(V,p)$ be a normal complex surface singularity and let
$\rho:\widetilde V\to V$ be a good resolution: $\widetilde V$ is smooth,
$\rho$ is proper and isomorphic away from $p$, and the full reduced
exceptional divisor $E=\sum_jE_j$ has simple normal crossings.
Its intersection matrix is negative definite.  Consequently, for a rational
divisor $Q$ on $\widetilde V$, there is a unique rational exceptional cycle
$\nu(Q)$ satisfying
\[
  \nu(Q)\cdot E_j=Q\cdot E_j\qquad\text{for every }j.
\]
Write the relative Zariski decomposition as
\[
 K_{\widetilde V}+E=P+N,
\]
where $N$ is an effective rational exceptional cycle, $P\cdot E_j\geq0$
for every $j$, and $P\cdot E_j=0$ whenever $E_j$ occurs in $N$.
Wahl's local logarithmic invariant \cite{Wahl1990} is
\begin{equation}\label{eq:beta-definition}
 \beta(V,p)=-\nu(P)^2.
\end{equation}
We also write $P_{\mathrm{loc}}=\nu(P)$ for this numerical exceptional cycle.
These are the relative decomposition and intersection conventions used by
Okuma \cite[Section~2, pp.~649--650]{Okuma2004}.

The value is unchanged on a higher good resolution and agrees with the
value on the minimal semigood resolution, where the exceptional divisor
has normal crossings, allowing irreducible components to have self-nodes.
Indeed, let $\pi:\widehat V\to\widetilde V$ be a further resolution,
isomorphic off $E$, with full reduced exceptional divisor $\widehat E$ over
$p$.  Log canonicity of the normal crossings pair gives
\[
 K_{\widehat V}+\widehat E
   =\pi^*(K_{\widetilde V}+E)+G,
 \qquad G\geq0,
\]
where $G$ is $\pi$-exceptional.  Hence its relative decomposition is
$\pi^*P+(\pi^*N+G)$: the first summand is relatively nef and is orthogonal
to every component of the second.  Uniqueness of the relative
decomposition applies.  Also $\nu(\pi^*P)=\pi^*\nu(P)$, as one checks
against both strict transforms and $\pi$-exceptional curves.  The square
therefore stays unchanged.  A common higher resolution compares any two
choices, including a semigood one.

\begin{lemma}[Projective computation of the local square]
\label{lem:projective-local-square}
Let $V$ be an integral normal projective complex surface, smooth outside
a point $p$, with Cartier canonical divisor.  Let $b:S\to V$ be a
projective birational morphism from a normal surface, isomorphic outside
$p$.  Suppose that the reduced inverse image of $p$ is a divisor $B$,
that $(S,B)$ is log canonical, and that $K_S+B$ is $b$-nef.  Then
\begin{equation}\label{eq:projective-local-square}
 \beta(V,p)=K_V^2-(K_S+B)^2.
\end{equation}
\end{lemma}

\begin{proof}
Choose a log resolution $\pi:Z\to S$, isomorphic outside $B$, and put
$\rho=b\pi$.  Let $E$ be the full reduced exceptional divisor of $\rho$.
With compatible canonical divisors, log canonicity gives
\[
 K_Z+E=\pi^*(K_S+B)+N,
\]
where $N$ is effective and $\pi$-exceptional.  The pullback on the right
is $\rho$-nef and has intersection zero with every component of $N$.
It is therefore the positive part $P$ of the relative decomposition.

The divisor $Z_P=P-\rho^*K_V$ is exceptional: the two summands agree
outside $\rho^{-1}(p)$.  Moreover, $\rho^*K_V$ has intersection zero with
every exceptional curve, so $Z_P=\nu(P)$.  The projection formula gives
\[
 P^2=K_V^2+Z_P^2=(K_S+B)^2.
\]
Together with \eqref{eq:beta-definition}, this proves the formula.
\end{proof}

\begin{lemma}[Persistence of log canonicity]\label{lem:lc-persistence}
In a holomorphic family of isolated hypersurface singularities in three
spatial variables with constant Milnor number over a connected smooth
complex curve, log
canonicity of the normal surface germ is constant.
\end{lemma}
\begin{proof}
By Varchenko's theorem \cite{Varchenko1982}, as recovered by
Steenbrink, the spectrum of an isolated hypersurface singularity
is constant in a $\mu$-constant deformation
\cite[Theorem~2.8]{Steenbrink1985}.  Its smallest number plus one is the
complex singularity index \cite[Theorem~2.11]{Steenbrink1985}, and the
log canonical threshold is the minimum of this index and one
\cite[Theorem~9.5]{Kollar1997Pairs}.  For a normal hypersurface in a smooth
ambient space, adjunction and inversion of adjunction identify log
canonicity with threshold one \cite{Kawakita2007}.  To apply the algebraic
inversion theorem to an individual analytic germ, replace its defining
function by a sufficiently high Taylor polynomial using
Lemma~\ref{lem:individual-jets}; analytic right equivalence preserves both
the threshold and log canonicity.  Thus log canonicity, and its failure,
is constant in the family.  This use of the spectrum concerns the
individual isolated function germs; it does not assume topological
triviality of the family.
\end{proof}

\section{Reduction of the numerical assertion to algebraic curves}
\label{sec:algebraic-reduction}

We now use the Wahl invariant $\beta=-P_{\mathrm{loc}}^2$ defined above.
The objective of this section is to transfer constancy of this scalar
invariant from algebraic curve families to arbitrary holomorphic disk
families.  The individual right-equivalences of
Lemma~\ref{lem:individual-jets} are sufficient; no simultaneous analytic
equivalence of the original family with a polynomial family is asserted.

Fix $\mu>0$ and $N\geq\mu+1$.  Let $V_N$ be the affine space of complex
polynomials in three variables of degree at most $N$, with zero constant
and linear terms, and put
\[
 U_\mu=\{g\in V_N:\dim_{\mathbb C}\mathbb C\{x\}/J(g)=\mu\}.
\]

\begin{lemma}[The precise jet locus]
\label{lem:mu-jet-locus}
The subset $U_\mu$ is locally closed algebraic.  On $U_\mu$, the function
$g\mapsto\beta(V(g),0)$ has finite image and is algebraically
constructible.
\end{lemma}

\begin{proof}
Work first in the finite-dimensional algebra
$A_\mu=\mathbb C[x_1,x_2,x_3]/\mathfrak m^{\mu+1}$.  The ideal generated
by the partials of $g$ is the image of a linear map
$A_\mu^{\oplus3}\to A_\mu$ whose matrix entries depend polynomially on
the coefficients of $g$.  The condition that its cokernel have
dimension $\mu$ is a locally closed rank condition.

This condition tests the full local length, including exclusion of
infinite length.  Indeed, for the local algebra
$A=\mathbb C\{x\}/J(g)$ with maximal ideal $\mathfrak n$, suppose
$\dim_{\mathbb C}A/\mathfrak n^{\mu+1}=\mu$.  The $\mu+1$ successive
layers $\mathfrak n^j/\mathfrak n^{j+1}$, $0\leq j\leq\mu$, cannot
all be nonzero.  If a layer vanishes, Nakayama's lemma applied to the
finitely generated ideal $\mathfrak n^j$ shows that this ideal is zero.
Consequently $A$ already has length $\mu$.  The converse follows from
the same filtration.  Polynomial, localized algebraic, and convergent
power-series computations give the same truncated quotient.  This
proves the first assertion.

For constructibility, let $B$ be any irreducible locally closed
subvariety of $U_\mu$; replace it first by a smooth dense open.
Consider its universal hypersurface with the origin section.  After
another dense open restriction on $B$, there is a Zariski neighborhood
of this section on which the relative singular locus is exactly the
section.  To see this, isolatedness in each fiber implies that the
section is an irreducible component of the relative critical locus:
a larger component containing it would have a positive-dimensional
fiber at a general section point, by the dimension inequality.
The other components meet the section over proper closed subsets of
$B$.  Remove those subsets from $B$, then remove those components
from the total space.  The resulting total hypersurface is
Cohen--Macaulay and regular off a subset of codimension two, hence
normal.  Its fibers have dimension two, so miracle flatness over the
smooth base gives flatness; it is smooth over $B$ off the section.

Take a projective resolution, isomorphic off the section, which
principalizes its ideal and has simple normal crossing reduced
exceptional divisor $E$.  Characteristic-zero resolution of
singularities and principalization provide this construction; see
\cite[Theorems~1.0.1--1.0.3]{Wlodarczyk2005}.
Generic smoothness, applied to the resolution space and every
intersection of components of $E$, permits a further dense open
restriction on $B$ such that the resolution is smooth over $B$ and
$E$ is a relative simple normal crossing divisor.  Nondominating
exceptional strata are discarded in making this restriction.
Since the resolution is projective and all of $E$ lies over the
section, $E$ and all intersections of its components are proper over
$B$.

The fibers are now good resolutions of the surface germs at the
origins.  Their weighted exceptional dual graphs are locally constant
in the complex topology, up to permutation of the vertices.
In fact smooth proper variation preserves the connected components
and genera of the exceptional curves, as well as their intersection
points and incidence.  Their self-intersections are degrees of the
normal line bundles obtained by restricting the divisor line bundles
on the resolution space.  These degrees are locally constant in
smooth proper curve families.  If a total divisor has several
components in a fiber, they are disjoint because that relative divisor
is smooth, so the same degree calculation applies to each component.

The weighted graph determines $\beta$.  It determines the exceptional
intersection matrix, and adjunction gives
$K\cdot E_i=2g(E_i)-2-E_i^2$ for each exceptional curve.  Thus it
determines all intersections of $K+E$ with the exceptional curves,
and hence its relative Zariski decomposition and the square of its
numerical exceptional representative.  An irreducible smooth complex
base is connected, so $\beta$ is constant on this dense open of $B$.
Noetherian induction on the remaining closed subsets yields a finite
stratification of $U_\mu$ on which $\beta$ is constant.
\end{proof}

\begin{proposition}[Scalar curve reduction]
\label{prop:scalar-curve-reduction}
Assume the following algebraic assertion: for every polynomial family
in three variables with regular coefficients on a smooth complex
algebraic curve $C$, with zero constant and linear terms and constant
positive Milnor number at the origin, the Wahl invariant at a marked
non-log-canonical fiber equals its value on the nearby general fibers.
Then $\beta(X_t,0)$ is constant near $t=0$ in every holomorphic
$\mu$-constant family whose central surface germ is non-log-canonical.
\end{proposition}

\begin{proof}
Take $N\geq\mu+1$ and consider the holomorphic coefficient map
$a:\Delta\to V_N$, $a(t)=j^Nf_t$.  By
Lemma~\ref{lem:individual-jets}, it takes values in $U_\mu$ and
preserves the analytic type of every individual surface germ, hence
its Wahl invariant.  Write $a_0=a(0)$ and $\beta_0=\beta(V(a_0),0)$.

If constancy fails on every smaller disk, $a_0$ lies in the closure
of the constructible set
\[
 Z_{\ne}=\{g\in U_\mu:\beta(V(g),0)\ne\beta_0\}.
\]
The closure here may be taken in the Zariski topology: an accumulating
sequence of coefficient vectors already forces this inclusion.
Choose an irreducible locally closed subset of $Z_{\ne}$ whose closure
in $U_\mu$ contains $a_0$.  After restricting to an affine neighborhood
of $a_0$, its closure is an irreducible affine variety $Z$ of positive
dimension, containing a dense open $Z^\circ\subset Z_{\ne}$.

There is an algebraic curve in $Z$ through $a_0$ whose general point
lies in $Z^\circ$.  For completeness, choose a finite Noether
normalization $\pi:Z\to\mathbb A^d$ and translate so that
$\pi(a_0)=0$.  The image of $Z\setminus Z^\circ$ is a proper closed
subset.  Choose a line through $0$ not contained in that image.
Its inverse image is cut out in $Z$ by $d-1$ equations, so every
irreducible component through $a_0$ has dimension at least one.
Finiteness of $\pi$ gives the opposite inequality and prevents a
one-dimensional component from mapping to a point.  Such a component
therefore dominates the line and has general point in $Z^\circ$.
Normalize it, choose a point over $a_0$, and remove the finitely many
other points needed to obtain the desired marked smooth curve $C$.
The inclusion into coefficient space supplies regular polynomial
coefficients on $C$.

Every member of this curve lies in $U_\mu$, its marked fiber is the
polynomial $a_0$ with non-log-canonical surface germ, and its general
Wahl invariant differs from $\beta_0$.  This contradicts the assumed
algebraic assertion.  Therefore the invariant is constant in the
original analytic family after shrinking.
\end{proof}

Appendix~\ref{app:local-numerical} gives a second curve comparison,
using irreducible components of $U_\mu$ directly; it does not require
constructibility of the scalar invariant.

\subsection{A projective family with the prescribed singularities}
\label{sec:compactification}

The following construction allows the local invariant to be compared by
ordinary intersection numbers on projective surfaces.  It preserves each
individual analytic right type; no simultaneous analytic coordinate
change is needed.

\begin{lemma}[Compactification with one singular point]
\label{lem:compactification}
Let $C$ be a smooth connected complex algebraic curve with marked point
$0$.  Let $g_c(x_1,x_2,x_3)$ be polynomials of bounded degree $l$, with
coefficients regular on $C$, satisfying $g_c(0)=0$ and
$d_xg_c(0)=0$.  Suppose the critical point at the origin is isolated and
has the same finite Milnor number $\mu$ for every $c\in C$.
After replacing $C$ by a Zariski neighborhood of $0$, there is a flat
projective hypersurface family $X\subset\mathbb P^3\times C$ of some
fixed degree $d$, with section $\Gamma$, such that each fibre is integral
and normal, smooth outside $\Gamma(c)$, and its germ at $\Gamma(c)$
admits a defining function analytically right-equivalent to $g_c$.
Furthermore,
\begin{align}
 K_{X_c}^2&=d(d-4)^2,\label{eq:compact-canonical-square}\\
 \chi_c(X_c\setminus\{\Gamma(c)\})
   &=d(d^2-4d+6)-\mu-1.\label{eq:compact-euler}
\end{align}
The total space $X$ is normal and Gorenstein.  These conclusions remain
valid after a finite change of smooth base curves.
\end{lemma}

\begin{proof}
Choose $d>l$ with $d\geq\mu+2$ and homogeneous coordinates
$[x_1:x_2:x_3:w]$ on $\mathbb P^3$.  Homogenize $g_c$ to degree $d$,
obtaining $G_c(x,w)$, and set
\[
 X_c=\{G_c(x,w)+Q_d(x)=0\},
 \qquad \Gamma(c)=[0:0:0:1],
\]
where $Q_d$ is a general homogeneous form of degree $d$ in the three
spatial coordinates, independent of $c$.

Since $Q_d\in\mathfrak m^{\mu+2}$, Lemma~\ref{lem:individual-jets}
shows, for each $c$ separately, that adding $Q_d$ preserves the analytic
right type of $g_c$.  This finite-determinacy and high-degree-tail
compactification is also the construction used in
\cite[Lemma~2.5]{Steenbrink1985}.

At $c=0$, the linear system spanned by $G_0$ and all forms $Q_d(x)$
has base locus exactly $\{[0:0:0:1]\}$.  The characteristic-zero
Bertini theorem \cite[Theorem~4.1]{KleimanBertini} therefore makes a
general member smooth elsewhere;
its nonzero coefficient of $G_0$ can be normalized to one.  One can
also see the required Bertini assertion directly.  On
$U=\mathbb P^3\setminus\{[0:0:0:1]\}$, the incidence variety
\[
 \{(z,Q):G_0(z)+Q(z)=0\}\subset U\times H^0(\mathbb P^2,\mathcal O(d))
\]
is an affine bundle over $U$, since evaluation of the forms is
surjective at every $z\in U$.  It is smooth.  Its generic fibre over
the coefficient space is regular, being obtained by localization, and
hence smooth over its characteristic-zero function field.  The
nonsmooth locus has constructible image missing the generic point,
so a nonempty open set of choices has smooth fibres.

There are no other singularities for complex parameters sufficiently
near $0$.  Indeed, otherwise projectivity gives a sequence of singular
points converging to the unique singular point of $X_0$.  They are then
in the chart $w=1$.  Conservation of the spatial-gradient number in a
small ball contradicts their existence, since the origin already
accounts for the full constant number $\mu$.  The parameters with an
additional singular point form a constructible subset of $C$.  A
constructible subset of an algebraic curve which misses an analytic
neighborhood of $0$ is finite; removing it gives the required Zariski
neighborhood.

Every resulting projective hypersurface has only one singular point.
It is reduced and irreducible: a repeated component or the intersection
of two distinct surface components would give a positive-dimensional
singular locus.  The hypersurface is Cohen--Macaulay and regular in
codimension one, hence normal.  Since the defining equation is nonzero
on every ambient fibre, it is a relative effective Cartier divisor and
the family is flat.  Outside the section it is smooth over $C$.
The total hypersurface is Gorenstein and regular in codimension one,
so it too is normal.

Adjunction gives $\omega_{X_c}=\mathcal O_{X_c}(d-4)$ and therefore
\eqref{eq:compact-canonical-square}.  A small global smoothing changes
only a Milnor neighborhood of the singular point: outside it the
family is a proper submersion with transverse boundary.  Replacing
the contractible singular neighborhood by its Milnor fibre increases
Euler characteristic by $\mu$.  A smooth degree-$d$ surface in
$\mathbb P^3$ has Euler characteristic $d(d^2-4d+6)$, as follows from
its tangent-normal exact sequence and the Gauss--Bonnet formula.
This proves \eqref{eq:compact-euler}.  Such a smoothing can, for example,
be obtained by subtracting $\varepsilon w^d$; near the isolated
critical point this is the usual Milnor smoothing, and compactness
preserves smoothness away from that neighborhood.

A change of base leaves the individual projective surfaces and their
germs unchanged.  Over a smooth new curve the relative hypersurface
description, flatness, smoothness off the section, and the same
codimension-one argument give all the remaining assertions again.
\end{proof}

\section{A semistable log model}\label{sec:log-model}

The algebraic comparison requires a model on which the exceptional curves
of the general surface form a horizontal boundary.  The following construction
keeps that boundary while making the special fibre reduced.  Its output is the
input for the support and adjunction arguments in the next section.

\begin{proposition}[The relative log model]\label{prop:log-model}
Let $C$ be a smooth connected complex algebraic curve with a marked point $0$,
and let $\pi:X\to C$ be a flat projective family of integral normal Gorenstein
surfaces.  Suppose that $X$ is normal and Gorenstein, that a section
$\Gamma\subset X$ contains every fibre singularity, and that $\pi$ is smooth
on $X\setminus\Gamma$.  Then, after shrinking $C$ about $0$ and making a
finite base change from a smooth curve, there is a projective birational
morphism
\[
 h:Y\longrightarrow X
\]
with the following properties.  Here $X$, $C$, $\Gamma$ and $0$ denote the
base-changed objects, and $H$ is the sum, with coefficient one, of all
$h$-exceptional prime divisors dominating $C$.
\begin{enumerate}
\item $Y$ is normal and $\mathbb Q$-factorial, and $h$ is an isomorphism
      over $X\setminus\Gamma$.
\item For every $c\in C$, the scheme fibre $Y_c$ is a reduced Cartier divisor,
      and $(Y,H+Y_c)$ is divisorial log terminal.
\item $K_Y+H$ is nef over $X$.  Every vertical $h$-exceptional prime divisor
      lies in $Y_0$ and maps to the point $\Gamma\cap X_0$.
\item For $c\ne0$, the surface $Y_c$ is integral and normal, and the divisor
      $H|_{Y_c}$ is reduced as a Weil divisor.  The fibre $Y_0$ has a unique component dominating
      $X_0$; its remaining components are projective surfaces over
      $\Gamma\cap X_0$.
\end{enumerate}
\end{proposition}

Here nefness over $X$ means nonnegative intersection with every complete curve
contracted by $h$.  A divisorial log terminal (dlt) pair is used in the usual
characteristic-zero sense; in particular its coefficient-one components are
normal and it has simple normal crossings at the generic points of their
intersection strata.  We construct the marked semistable model by
Proposition~\ref{prop:marked-monomial-reduction}, which makes explicit the
dimension-three monomial reduction of
\cite[Theorems~7.17 and~7.19]{KM1998}, and then apply the semistable minimal
model theorem \cite[Theorem~7.9]{KM1998}.  In its terminology, a dlt morphism
to a smooth curve requires the fibre to be included in the boundary at every
parameter \cite[Definition~7.1]{KM1998}.

\begin{proof}
First choose a projective log resolution $g:Z\to X$ which principalizes the
ideal of $\Gamma$ and resolves the union with $X_0$.  We choose it to be an
isomorphism over $X\setminus\Gamma$: that open set is smooth and its special
fibre is already smooth.  Thus $Z$ is smooth, $g^{-1}(\Gamma)$ has divisorial
support, and the union of this support with the support of $Z_0$ has simple
normal crossings.  All exceptional divisors of $g$ lie over $\Gamma$.
By generic smoothness for $Z\to C$ and for the finitely many strata of its
exceptional divisor, we may delete finitely many base points other than $0$
so that $Z\to C$ is smooth away from $0$ and every fibre support together with
the exceptional divisor has simple normal crossings.  After the same
shrinking, every vertical exceptional divisor lies over $0$.

Let $H_Z$ be the reduced horizontal exceptional divisor of $g$, and write
$Z_0=\sum a_jZ_j$.  The preceding shrinking makes
$(Z_0)_{\mathrm{red}}+H_Z$ simple normal crossing, with $H_Z$ relatively
simple normal crossing over $C\setminus\{0\}$.  Choose a positive integer
$b$ divisible by every $a_j$, and a parameter $\tau$ at $0$ regular with no
other zero on a Zariski neighbourhood.  The cover $s^b=\tau$, with its smooth
normalization, is finite, ramified to order $b$ over $0$, and \'etale
elsewhere.  We again write $X$, $C$ and $\Gamma$ for the pulled-back objects.

The pulled-back target $X$ is still normal and Gorenstein.  Indeed the chosen
curve cover is locally $s^b=\tau$, so the target is locally the hypersurface
$A[s]/(s^b-\tau)$ over a Gorenstein ring $A$.  It is therefore Gorenstein and
satisfies Serre's condition $S_2$.  Off the section it is smooth over the new
smooth curve, while the section has codimension two.  Thus it satisfies
$R_1$ and is normal.  Its fibres remain integral, so the flat total space over
the connected covering curve is integral.

Apply Proposition~\ref{prop:marked-monomial-reduction} to $Z$ with marking
$H_Z$ and the chosen cover.  It gives a smooth variety $V$, projective over
the normalized pullback of $Z$, whose special fibre is reduced and simple
normal crossing together with the transformed horizontal marking.
The construction is an isomorphism on the normalized pullback away from $0$,
where the family and marking already have these properties.  Over
$X\setminus\Gamma$, the original morphism was smooth and the marking was
empty, so the same assertion shows that $V\to X$ is an isomorphism there.
Thus all exceptional divisors still lie over $\Gamma$, and every
vertical exceptional divisor lies over $0$.

Let $H_V$ be the reduced sum of the horizontal exceptional prime divisors on
$V$.  Since the modification is an isomorphism away from $0$, $H_V$ is
precisely the transformed reduced marking $H_Z$.  For every $c$, the pair
$(V,H_V+V_c)$ is simple normal crossing.
Consequently $(V,H_V)\to C$ is a dlt morphism.  The relative minimal model
program of \cite[Theorem~7.9]{KM1998} applies to the threefold pair $(V,H_V)$
over the flat target $X\to C$.  All its contractions and flips are projective
over $X$, it terminates, and the resulting morphism $(Y,H)\to C$ is again
dlt.  Its alternative of a Mori fibre space cannot occur: such a space would
factor the dominant morphism from the threefold to $X$ through a variety of
dimension less than three.  We therefore obtain $K_Y+H$ nef over $X$.

The program consists of divisorial contractions and flips, so it extracts
no divisors.  Every surviving horizontal exceptional prime consequently
retains coefficient one in $H$, and every vertical exceptional prime still
lies over $0$.  Each step is an isomorphism over $X\setminus\Gamma$, since
there were no relative curves to contract there.  The usual preservation of
$\mathbb Q$-factoriality in this program gives the asserted
$\mathbb Q$-factorial model.

For completeness, the scheme-theoretic fibre is precisely the boundary fibre
used above.  The pullback of a local parameter at $c$ is a nonzerodivisor on
the integral normal variety $Y$, so $Y_c$ is Cartier.  Its order at every prime
component is one: the corresponding divisorial valuation occurred on $V$,
where the fibre was reduced, and the program did not extract divisors.
A principal divisor with these orders on a normal variety defines a reduced
subscheme.  Locally, if $u^n\in(t)$, the height-one valuations show that
$u/t$ is regular at every height-one prime; normality then gives $u/t\in
\mathcal O_Y$.  Thus $(t)$ is radical.

For $c\ne0$ there are no vertical exceptional divisors, and $X_c$ is integral.
Its strict transform is therefore the only component of $Y_c$.  This component
is normal by the dlt property.  At the generic point of each curve in
$H\cap Y_c$, the pair is simple normal crossing; hence the Cartier fibre cuts
$H$ with multiplicity one, giving the reduced Weil divisor $H|_{Y_c}$.  Over $0$
there is again just one strict transform of $X_0$; all other components are
exceptional.  They map into $\Gamma\cap X_0$ and are projective, because $h$
is projective.  This proves all the assertions.
\end{proof}

The proposition identifies $H$ as the full surviving horizontal divisorial
boundary.  It does not yet identify the whole set $h^{-1}(\Gamma)$ with the
boundary, nor identify the different in adjunction on each fibre component.
Those are separate consequences of nefness, non-log-canonicity and the support
argument proved next.

\subsection{The full exceptional boundary}
\label{sec:full-boundary}

The log model must remember every point over the singular section.  This
section proves that relative nefness forces that property in the
non-log-canonical case.  It also identifies precisely the boundary entering
surface adjunction.  The existence of the model is an input to this section.

Let $C$ be a smooth connected complex curve with distinguished point $0$,
let $\sigma:X\to C$ be a flat family of normal Gorenstein surfaces, and let
$\Gamma\subset X$ be a section.  Assume that $X$ is normal and Gorenstein.
Consider a projective birational morphism
\[
 h:Y\longrightarrow X,\qquad q=\sigma\circ h,
\]
where $Y$ is a normal $\mathbb Q$-factorial threefold and $h$ is an
isomorphism over $X\setminus\Gamma$.  Write $H$ for the reduced sum of the
$h$-exceptional prime divisors dominating $C$.  Assume that all vertical
$h$-exceptional prime divisors lie over $0$, that every fiber $Y_c$ is a
reduced Cartier divisor, and that
\[
 Y_0=S_0+\sum_a S_a,
\]
where $S_0$ maps birationally onto $X_0$ and every $S_a$ maps to the point
$\Gamma\cap X_0$.  We assume that $Y_c$ is irreducible for $c\ne0$, that
$(Y,H+Y_c)$ is dlt for all $c$ under consideration, and that
\[
 L=K_Y+H
\]
is $h$-nef.  We use compatible canonical divisors and put
\[
 A=L-h^*K_X,\qquad D=Y_0+H.
\]
The assumptions on the fibers may equally be imposed after shrinking $C$.
For a component $S$ of $Y_c$, let $B_S$ be the reduced union of the curves
in which $S$ meets the other components of $H+Y_c$.

\begin{lemma}[Surface adjunction]
\label{lem:surface-boundary}
Every component $S$ of $Y_c$ is normal.  Its different is exactly $B_S$:
\[
 (K_Y+H+Y_c)|_S\sim_{\mathbb Q}K_S+B_S.
\]
In particular $(S,B_S)$ is dlt and
\[
 L|_S\sim_{\mathbb Q}K_S+B_S.
\]
The latter divisor is nef on every curve contracted by $h|_S$.  For an
extra vertical component $S_a$, it is nef on the projective surface $S_a$.
\end{lemma}

\begin{proof}
Dlt adjunction gives normality of $S$ and a dlt pair
$(S,\operatorname{Diff}_S(H+Y_c-S))$; see
\cite[Lemma~10.2]{Fujino2011}.  We check the coefficients of the different
at every prime curve of $S$.

If a prime curve lies in another component of $H+Y_c$, its generic point
is an lc stratum.  The dlt pair is simple normal crossing there, so ordinary
adjunction gives coefficient one.  Here we use the description of dlt lc
centers as boundary strata and the defining simple-normal-crossing property
at their generic points; see
\cite[Definition~11.5 and (11.10.4)]{Kollar2023Families}.

At the generic point $\eta$ of any other prime curve of $S$, no other
component of $H+Y_c$ is present.  Thus $S$ equals the Cartier divisor
$Y_c$ near $\eta$.  The local ring $\mathcal O_{S,\eta}$ is a regular
one-dimensional local ring, since $S$ is normal.  If $t$ is a local
equation of the fiber, the two-dimensional local ring
$R=\mathcal O_{Y,\eta}$ satisfies
\[
 R/(t)=\mathcal O_{S,\eta}.
\]
Consequently its maximal ideal is generated by $t$ and one lift of a
generator of the maximal ideal of $R/(t)$.  The embedding dimension of $R$
is at most two, and $R$ is regular.  Ordinary adjunction therefore gives
coefficient zero at this curve.  This proves equality of the different
and $B_S$ as divisors; no other prime-curve contribution remains.

The line bundle of $Y_c=q^*(c)$ restricts trivially to $S$, giving the
formula with $L$.  Its degree on a curve contracted by $h|_S$ is the degree
of $L$ on that curve and is nonnegative.  Finally, $S_a$ is projective
because it is a closed subvariety of the projective fiber of $h$ over
$\Gamma\cap X_0$, and every curve on $S_a$ is contracted by $h$.
\end{proof}

\begin{proposition}[Support over a non-log-canonical section]
\label{prop:full-exceptional-boundary}
Assume in addition that the surface $X_c$ is not log canonical at
$\Gamma\cap X_c$ for general $c\in C$.  Then $A\le0$, some horizontal coefficient of $A$ is
strictly negative, and
\[
 h^{-1}(\Gamma)=\operatorname{Supp}(-A)
              =\operatorname{Supp}(H)\cup\bigcup_a S_a
\]
as sets.  In particular $B_{S_0}$ is the full reduced exceptional
preimage of $\Gamma\cap X_0$ on $S_0$, and $B_{Y_c}$ has the same property
on every $Y_c$ with $c\ne0$.  Thus
\[
 S_0\setminus\operatorname{Supp}(B_{S_0})
       \simeq X_0\setminus\Gamma,
 \qquad
 Y_c\setminus\operatorname{Supp}(B_{Y_c})
       \simeq X_c\setminus\Gamma\quad(c\ne0).
\]
\end{proposition}

\begin{proof}
The divisor $A$ is exceptional, since $h$ is an isomorphism away from
$\Gamma$, and is $h$-nef, since a pullback from $X$ has degree zero on
contracted curves.  Apply the negativity lemma to $-A$: its pushforward
is zero, so $-A$ is effective
\cite[Lemma~4.16(1)]{Fujino2011}.

Suppose every horizontal coefficient of $A$ were zero.  The support of
$A$ would then lie over $0$.  For general $c\ne0$, its restriction to
$Y_c$ would be zero.  By Lemma~\ref{lem:surface-boundary} and Cartier-fiber
adjunction on $X$, with compatible canonical divisors, this gives a
crepant comparison
\[
 K_{Y_c}+B_{Y_c}=h_c^*K_{X_c},\qquad h_c=h|_{Y_c}.
\]
The pair on the left is log canonical.  Computing discrepancies on a
common resolution now shows that $X_c$ is log canonical: the discrepancy
of a divisor exceptional over $Y_c$ is unchanged, and a boundary component
on $Y_c$ has discrepancy $-1$ over $X_c$.  This contradicts the hypothesis.
Thus at least one horizontal exceptional prime has negative coefficient.

Such a prime maps onto $\Gamma$: its image is closed by properness, is
contained in $\Gamma$, and dominates $C$.  Therefore every fiber of $h$
over $\Gamma$ meets $\operatorname{Supp}(-A)$.  The support statement in
the negativity lemma says that each fiber of a proper birational morphism
between normal varieties either lies in the support of an effective
anti-nef divisor or is disjoint from it
\cite[Lemma~4.16(2)]{Fujino2011}.  Applying it to $-A$ proves
$h^{-1}(\Gamma)\subset\operatorname{Supp}(-A)$; the reverse inclusion
holds because $A$ is exceptional over $\Gamma$.

Every prime in $\operatorname{Supp}(-A)$ is either a component of $H$ or
one of the $S_a$, by the assumptions on the exceptional primes.
Conversely these divisors map into $\Gamma$ and hence lie in the support.
Intersect this equality with $S_0$.  None of the listed divisors equals
$S_0$, and their intersections are precisely the boundary curves defining
$B_{S_0}$; dlt strata have the expected codimension.  Hence there are no
additional exceptional points or curves on $S_0$ outside this boundary.
The same argument on $Y_c$, $c\ne0$, uses only $H$.  The stated
isomorphisms follow from the assumed isomorphism away from $\Gamma$.
\end{proof}

\begin{remark}
Lemma~\ref{lem:surface-boundary} identifies a divisor on a surface.  It
does not assert that further adjunction to a boundary curve has zero
fractional different, nor that the threefold is smooth at every closed
boundary point.  Those are distinct local questions.
\end{remark}

\section{The nearby Euler characteristic along the boundary}
\label{sec:nearby-euler}

We now compare a nearby open fiber of the log model with the open
surfaces in its special fiber.  The essential local assertion is that
boundary intersections contribute zero, including the points where the
threefold is singular.  We first prove that assertion independently of
the global construction.

For a complex algebraic variety, write $\chi_c$ for the Euler
characteristic with compact supports.  On a complex algebraic variety it
equals the ordinary Euler characteristic.  If $f:Z\to\mathbb C$ is
regular, the \emph{nearby Euler function} $\psi_f$ is the unshifted
nearby-cycle operation on constructible functions.  In particular,
$\psi_f(\mathbf1_V)(z)$ is the Euler characteristic of the local Milnor
fiber of $f|_V$ when $V$ is closed.  Extend this definition by additivity
to constructible functions.  Equivalently, take stalk Euler
characteristics of the usual, unshifted sheaf nearby cycles.  This
operation depends only on the restriction away from $f^{-1}(0)$ and
commutes with proper pushforward.  For a finite map $\rho$, pushforward
is the sum over the finite set $\rho^{-1}(z)$; for a general proper map it
is integration with respect to $\chi_c$.  These conventions and the
proper compatibility are recalled in
\cite[Definition~3.2]{SchuermannNearby} and
\cite[Proposition~4.19(1)]{MaximSchuermann}.

\begin{lemma}[Boundary covers in dimension three]
\label{lem:boundary-cover}
Let $Z$ be a normal complex algebraic threefold and let
$B=\sum_{j=1}^n B_j$ be a reduced divisor with distinct prime components.
Assume that $(Z,B)$ is dlt and that every $B_j$ is $\mathbb Q$-Cartier.
If $z$ lies on at least two of the $B_j$, then, after replacing $Z$ by a
neighborhood $U$ of $z$, there is a finite surjective map
$\rho:\widetilde U\to U$ from a normal variety such that:
\begin{enumerate}
\item $\rho$ is \'etale in codimension one and over
      $U\setminus B$, and every total pullback $\rho^*B_j$ is a reduced
      Cartier divisor;
\item $(\widetilde U,\rho^*B)$ is dlt and is simple normal crossing at
      every point of $\rho^{-1}(z)$.
\end{enumerate}
\end{lemma}

\begin{proof}
Shrink $U$ so that the only boundary components present are those
through $z$, and choose regular functions $g_j$ with
$\operatorname{div}(g_j)=m_jB_j$, where $m_j$ is a positive integer.
Normalize $U$ in a field extension obtained by adjoining roots
$a_j^{m_j}=g_j$ for all $j$.  Normalization is finite here, and taking
one field extension gives a finite surjective map of some degree $d$.
This is the usual local boundary-Cartier cover; see also
\cite[Proposition~11.25]{Kollar2023Families}.

At a codimension-one point, write each $g_j$ as a unit times the
$m_j$-th power of a local equation, or as a unit if $B_j$ is absent.
Adjoining its root is therefore unramified there.  The same argument
shows that the cover is \'etale wherever all the $B_j$ are Cartier,
in particular over $U\setminus B$ and over the simple normal crossing
locus of $(U,B)$.  Consequently
$\operatorname{div}(a_j)=\rho^*B_j$ has coefficient one on every prime
in its support.  Thus the total pullback is reduced and Cartier.

Put $\widetilde B=\rho^*B$.  The codimension-one ramification formula
gives
\[
 K_{\widetilde U}+\widetilde B
   =\rho^*(K_U+B)
\]
with compatible canonical divisors.  More generally, if a divisorial
valuation $F$ upstairs restricts to $e$ times a normalized divisorial
valuation $E$ downstairs, then
\begin{equation}
\label{eq:finite-log-discrepancy}
 A(F,\widetilde U,\widetilde B)=eA(E,U,B),
\end{equation}
where $A$ denotes log discrepancy.  To see the formula, choose a
birational model on which $E$ is a divisor and normalize it in the
upper function field.  At the generic point of $F$, the ramification
divisor has coefficient $e-1$.  Hence the discrepancy coefficients
satisfy $a_F=e a_E+(e-1)$, giving
\eqref{eq:finite-log-discrepancy}.  Restriction of a divisorial
valuation through a finite extension is divisorial, since its value
group has finite index and its residue field has the same
transcendence degree.  This is the finite-cover discrepancy formula
of \cite[Proposition~5.20]{KM1998}; see also
\cite[11.22--11.23]{Kollar2023Families}.

It follows that the upper pair is log canonical and that every log
canonical center upstairs maps onto a log canonical center downstairs.
Each lower center meets the simple normal crossing locus because the
lower pair is dlt.  The cover is \'etale there, so each upper center
also meets the simple normal crossing locus.  This is the dlt
criterion.  We use the standard facts that the log canonical centers
of a dlt pair are the normal strata of its coefficient-one boundary,
with the expected codimensions, and that intersections of log
canonical centers are unions of such centers; see
\cite[Proposition~9.2 in the author draft]{FujinoDlt} and
\cite[Definition~11.5 and~11.10]{Kollar2023Families}.

Fix $y\in\rho^{-1}(z)$.  There is at least one upper boundary prime
through $y$ over each lower boundary prime through $z$.  Indeed each
root $a_j$ vanishes at $y$, so its Cartier divisor contains $y$.
Primes over distinct lower primes are distinct, since a finite map
preserves their dimensions.  Thus at least two upper primes meet at
$y$.  If at least three meet, the dlt stratum theorem makes $y$ a
zero-dimensional log canonical center.  The dlt criterion then says
that the upper pair is simple normal crossing at $y$.

Suppose instead that precisely two upper boundary primes $T_1,T_2$
meet at $y$.  They lie over distinct lower primes; hence each agrees
near $y$ with its entire reduced Cartier pullback and is itself
Cartier there.  The surface $T_1$ is normal by dlt adjunction, and a
normal surface is Cohen--Macaulay.  Therefore the Cartier curve
$C=T_2|_{T_1}$ has no embedded points.  It is generically reduced,
because the pair is simple normal crossing at the generic points of
its curve strata, so it is reduced.  Each irreducible component of
$C$ is a normal curve, hence smooth.  Two such components cannot meet
at $y$: their intersection would make $y$ a log canonical center,
forcing simple normal crossings there, whereas a simple normal
crossing threefold pair with only two boundary branches has no
zero-dimensional log canonical center.  Consequently $C$ is a smooth
curve near $y$.

Let $R=\mathcal O_{\widetilde U,y}$ and let $a,b$ define $T_1,T_2$.
Since $R/(a,b)$ is a regular local ring of dimension one,
\[
 \operatorname{edim}R
 \le \operatorname{edim}(R/(a,b))+2=3=\dim R.
\]
Thus $R$ is regular.  Moreover $a,b$ have independent images in its
cotangent space, so they are part of a regular system of parameters.
The two boundary components are transverse.  This proves the claimed
simple normal crossing assertion in every case.
\end{proof}

\begin{lemma}[Vanishing at boundary intersections]
\label{lem:boundary-nearby-zero}
Under the hypotheses of Lemma~\ref{lem:boundary-cover}, suppose that a
regular function $f$ has reduced zero divisor
$V=\sum_{j\in I}B_j$, where $I$ is a nonempty subset of
$\{1,\ldots,n\}$.  Put $H=\sum_{j\notin I}B_j$.
At every point $z\in V$ lying on at least two components of $B$,
\[
 \psi_f(\mathbf1_{Z\setminus H})(z)=0.
\]
\end{lemma}

\begin{proof}
Use the finite cover of Lemma~\ref{lem:boundary-cover}, and let $d$
denote its degree.  At each point
$y$ over $z$, choose analytic simple normal crossing coordinates.  As
the cover is unramified in codimension one, the vertical
multiplicities remain one.  After absorbing a unit into one
coordinate, the pulled-back function and horizontal divisor are
\[
 f\circ\rho=z_1\cdots z_r,
 \qquad
 \rho^{-1}H=\{z_{r+1}\cdots z_{r+h}=0\},
 \qquad r\ge1,\quad r+h\ge2.
\]
When $h=0$, the horizontal divisor is empty.  For $r\ge2$, the local
Milnor fiber of the coordinate product has the homotopy type of
$(S^1)^{r-1}$ and therefore Euler characteristic zero.  The same is
true after restricting to any intersection of horizontal coordinate
hyperplanes.  Inclusion--exclusion therefore gives zero for the
complement of the horizontal boundary.  For $r=1$, each such
restricted Milnor fiber is a polydisc, with Euler characteristic
one.  Here $h\ge1$, so inclusion--exclusion gives
$\sum_{J\subseteq\{1,\ldots,h\}}(-1)^{|J|}=(1-1)^h=0$.
These computations may be made on small adapted polydiscs, or on
Milnor balls, with the same local nearby Euler values.

Let $\widetilde\alpha=\mathbf1_{\widetilde U\setminus\rho^{-1}H}$.
Over $U\setminus V$, its finite pushforward is
$d\mathbf1_{U\setminus H}$: away from $H$ the cover is \'etale of
degree $d$, and on $H$ all the terms are zero.  Since nearby cycles
only use the restriction off $V$, proper compatibility gives
\[
 d\,\psi_f(\mathbf1_{U\setminus H})(z)
  =\sum_{y\in\rho^{-1}(z)}
    \psi_{f\circ\rho}(\widetilde\alpha)(y)=0.
\]
This is an identity of constructible Euler functions; no triviality
of the covering local system is needed.  Since $d>0$, the result
follows.
\end{proof}

\begin{proposition}[Euler comparison for the log model]
\label{prop:nearby-euler-comparison}
Let $q:Y\to C$ be a projective flat morphism from a normal complex
threefold to a smooth pointed curve $(C,0)$.  Let $H$ be a reduced
divisor with horizontal prime components.  Assume that the Cartier
fiber $Y_0=\sum S$ is reduced, that $(Y,Y_0+H)$ is dlt, and that its
boundary components are $\mathbb Q$-Cartier.  Assume also that
$Y_t\setminus H$ is smooth for all sufficiently small $t\ne0$.
For each component $S$ of $Y_0$, let $B_S$ be the reduced union of its
intersections with the other components of $Y_0+H$, and put
$S^\circ=S\setminus B_S$.  If each local smoothing of an isolated
singularity in $S^\circ$ has Milnor fiber Euler characteristic at
least one, then
\begin{equation}
\label{eq:nearby-euler-comparison}
 \chi_c(Y_t\setminus H)
   =\sum_{S\subset Y_0}\chi_c(S^\circ)+\delta,
 \qquad \delta\ge0,
\end{equation}
for all sufficiently small $t\ne0$.
\end{proposition}

\begin{proof}
Write $\lambda=\psi_q(\mathbf1_{Y\setminus H})$ using a local
coordinate on $C$ at $0$.  Proper compatibility gives
\[
 \chi_c(Y_t\setminus H)=\int_{Y_0}\lambda\,d\chi_c.
\]
At every point lying on at least two components of $Y_0+H$,
Lemma~\ref{lem:boundary-nearby-zero} gives $\lambda=0$.
Every remaining point lies in exactly one $S^\circ$.  The surface
$S$ is normal by dlt adjunction, so its singular points form a finite
set.  At a smooth point of $S^\circ$, the local reduced Cartier fiber
is precisely $S$.  Its regular local ring is a quotient of the
ambient local ring by the base parameter, so the ambient ring is
regular and that parameter is part of a regular system of
parameters.  Thus $q$ is smooth there and $\lambda=1$.

At a singular point $p\in S^\circ$, a neighborhood misses $H$ and all
other special-fiber components.  Flatness and the hypothesis on the
nearby open fibers therefore give a smoothing of the normal isolated
surface singularity $(S,p)$.  If $M_p$ is its Milnor fiber, then
$\lambda(p)=\chi(M_p)\ge1$ by the assumed interior smoothing
statement.  Integrating the disjoint open surfaces and these finitely
many corrections yields \eqref{eq:nearby-euler-comparison}, with
\[
 \delta=\sum_{S\subset Y_0}
        \sum_{p\in\operatorname{Sing}(S)\cap S^\circ}
             \bigl(\chi(M_p)-1\bigr)\ge0.
\]
\end{proof}

In the projective log model $h:Y\to X$, the special-fiber components
are $S_0$ and the extra surfaces $S_a$.  The preceding proposition
applies once the interior smoothing inequality is supplied.  In
particular, the boundary calculation does not require the different
obtained by further adjunction to a curve to have integral
coefficients.

\subsection{Interior Milnor fibers and the surface inequality}
\label{sec:surface-estimates}

We need two signs in the nearby Euler comparison.  The local smoothing of an
interior singularity contributes at least the value~$1$ contributed by a smooth
point.  On the other hand, the complement of a reduced boundary in an extra
vertical surface has nonnegative Euler characteristic because its log canonical
divisor is nef.  We prove these statements separately.

We use ordinary Euler characteristic for compact local Milnor fibers.
All Betti numbers below have rational coefficients.

\begin{lemma}[The interior correction]
\label{lem:interior-correction}
Let $q:(\mathcal V,x)\to(\mathbb C,0)$ be a flat holomorphic germ whose
scheme central fiber $(S,x)$ is a normal surface with an isolated singularity
and whose sufficiently small nonzero fibers are smooth.  If $F_x$ is a
sufficiently small Milnor fiber of $q$, then
\[
  \chi(F_x)\geq 1.
\]
At a smooth point of a reduced Cartier central surface in a normal threefold,
the local nearby Euler characteristic is~$1$.
\end{lemma}

\begin{proof}
Greuel--Steenbrink's theorem \cite{GreuelSteenbrink1983} gives
$b_1(F_x)=0$ for a smoothing of a normal isolated singularity; no
Gorenstein hypothesis is required.  We use the formulation in
\cite[Theorem~2.18]{Greuel2019}.
Choose a Milnor representative with smooth transverse boundary.  Its interior
is a smooth Stein surface, being a closed analytic submanifold of a small
open ball slice.  The boundary collar makes this interior homotopy equivalent
to the compact Milnor representative.  By the Andreotti--Frankel
dimension bound \cite{AndreottiFrankel1959}, a complex Stein surface has
the homotopy type of a CW complex of real dimension at most two; see also
\cite[Theorem~1 and the following remark~(2)]{HillNacinovich1993}.
Consequently $b_j(F_x)=0$ for $j>2$.  The Milnor fiber is nonempty, so
\[
  \chi(F_x)=b_0(F_x)+b_2(F_x)\geq1.
\]

For the last assertion, let $t$ denote the pulled-back base coordinate.
At the point in question, the quotient of the threefold local ring by $(t)$
is regular of dimension two.  Lifting generators of its maximal ideal and
adjoining $t$ shows that the total local ring has embedding dimension at most
three, hence is regular.  The same dimension count shows that $t$ is a regular
coordinate.  Thus $q$ is smooth there and its small local fiber is contractible.
\end{proof}

\begin{lemma}[The surface inequality]
\label{lem:surface-bmy}
Let $S$ be a normal projective surface over $\mathbb C$, and let $B$ be a
reduced effective Weil divisor, possibly empty.  Suppose that $(S,B)$ is log
canonical and that the $\mathbb Q$-Cartier divisor $K_S+B$ is nef.  Set
$S^\circ=S\setminus\operatorname{Supp}B$.  Then
\[
  0\leq (K_S+B)^2
    \leq 3e_{\mathrm{orb}}(S,B)
    \leq 3\chi_c(S^\circ),
\]
where $e_{\mathrm{orb}}$ denotes Langer's orbifold Euler number.
\end{lemma}

\begin{proof}
Nefness on the projective surface implies pseudoeffectivity and nonnegative
self-intersection.  The latter can also be checked after pulling back to a
projective resolution, where the divisor is still nef and has the same square.
The pseudoeffective form of Langer's inequality gives
$(K_S+B)^2\leq3e_{\mathrm{orb}}(S,B)$
\cite[Corollary~5.2]{Langer2003}.

It remains to compare the two Euler numbers.  By Langer's definition,
smooth points of the boundary have weight zero.  At every singular point of
the pair lying on that boundary the local orbifold number is also zero, since the boundary
has a coefficient-one branch and the pair is log canonical
\cite[Theorem~7.6]{Langer2003}.  For
$Z=\operatorname{Sing}S\setminus\operatorname{Supp}B$, the weighted formula
therefore reads
\[
  e_{\mathrm{orb}}(S,B)
    =\chi_c(S^\circ)
      +\sum_{x\in Z}\bigl(e_{\mathrm{orb}}(x;S,0)-1\bigr).
\]
Every term in the sum is nonpositive by the local upper bound
\cite[Definition~3.4 and Corollary~7.8]{Langer2003}.  This proves the last
inequality, including all singular-point corrections.
\end{proof}

For application, write the reduced special fiber of the model as
\[
  Y_0=S_0+\sum_a S_a,
\]
where $S_0$ dominates the original special surface and each $S_a$ is an extra
component contracted to its marked point.  Let $B_S$ be the reduced boundary
on a component $S$ supplied by adjunction, and put
$S^\circ=S\setminus\operatorname{Supp}B_S$.
By Proposition~\ref{prop:log-model} and Lemma~\ref{lem:surface-boundary},
every $S_a$ is projective, $(S_a,B_{S_a})$ is dlt and hence log canonical, and its log
canonical divisor is nef: all its curves are contracted by the morphism over
which the threefold log divisor is nef.  Lemma~\ref{lem:surface-bmy} consequently
applies to each extra component.  It need not apply to $S_0$.

Proposition~\ref{prop:full-exceptional-boundary} identifies the nearby
complements $Y_c\setminus H$ with the smooth punctured fibers
$X_c\setminus\Gamma$ for $c\ne0$.  Thus an interior singular point
of any $S_a$ is smoothed by the family.  Indeed the central fiber is locally that component alone, is
normal and reduced Cartier, and the nearby fibers there are smooth.  Its
contribution in Lemma~\ref{lem:interior-correction} exceeds the smooth-point
value by the nonnegative number $\chi(F_x)-1$.

Lemma~\ref{lem:boundary-nearby-zero} supplies the vanishing of nearby
Euler values at intersections of boundary components, and
Proposition~\ref{prop:nearby-euler-comparison} combines it with these
interior corrections.  Theorem~\ref{thm:beta-constant} will use that
comparison and the two signs above to force every extra logarithmic square
to vanish.  The orbifold boundary calculation in the surface inequality
concerns a different invariant and does not supply the nearby-Euler
vanishing.

\section{Specializing the logarithmic square}
\label{sec:square-specialization}

We now isolate the numerical consequence of the extra-component
vanishing.  All intersections in this section are ordinary projective
intersection numbers; the conversion to the local invariant is exactly
Lemma~\ref{lem:projective-local-square}.

\begin{lemma}[Square specialization]
\label{lem:square-specialization}
Let $q:Y\to C$ be a flat projective morphism to a smooth connected
complex curve, with pure two-dimensional fibres.  Let $L$ be a
rational Cartier divisor on $Y$.  If
$[Y_0]=\sum_j m_j[S_j]$ is the cycle of a fibre, then for every $c\in C$
\[
 (L|_{Y_c})^2=\sum_j m_j(L|_{S_j})^2.
\]
In particular, the multiplicities on the right are one when $Y_0$ is
reduced.
\end{lemma}

\begin{proof}
Choose $r>0$ such that $rL$ is Cartier.  Flatness and properness imply
that $\chi(Y_c,\mathcal O_Y(nrL)|_{Y_c})$ is constant in $c$ for every
integer $n$ \cite[Lemma~36.32.2]{StacksProject}.  This is a polynomial
in $n$ of degree at most two, whose quadratic coefficient is
$r^2(L|_{Y_c})^2/2$.  Its leading term is the sum of the leading terms
on the two-dimensional components, weighted by their generic lengths
\cite[Lemmas~33.45.1, 33.45.2 and 33.45.6]{StacksProject}.  Comparing
these coefficients and dividing by $r^2$ proves the assertion.
\end{proof}

\begin{proposition}[The numerical conclusion of the log model]
\label{prop:beta-from-log-model}
Let $X\to C$, with section $\Gamma$, be the projective family of
Lemma~\ref{lem:compactification}.  Let $h:Y\to X$ be a projective
birational morphism from a normal integral threefold, isomorphic off
$\Gamma$.  Suppose $Y_0=S_0+\sum_a S_a$ is reduced Cartier, $S_0$
maps birationally to $X_0$, and the $S_a$ map to $\Gamma(0)$.  Suppose
that, after shrinking $C$, the other fibres $Y_c$ are normal and map
birationally to $X_c$.

Let $H$ be the reduced horizontal exceptional boundary and set
$L=K_Y+H$, assumed rational Cartier and $h$-nef.  For each component
$S$ of $Y_0$, let $B_S$ be the reduced divisor of its intersections
with the other components of $Y_0+H$.  Assume that $(S,B_S)$ is log
canonical and
\[
 L|_S\sim_{\mathbb Q} K_S+B_S.
\]
For $S_0$, require that $B_{S_0}$ is the full reduced inverse image of
$\Gamma(0)$.  For $c\ne0$, require likewise that the full reduced
inverse image $B_c$ of $\Gamma(c)$ is a divisor, that $(Y_c,B_c)$ is
log canonical, and that
$L|_{Y_c}\sim_{\mathbb Q}K_{Y_c}+B_c$.
If
\[
 (K_{S_a}+B_{S_a})^2=0\qquad\text{for every }a,
\]
then $\beta(X_c,\Gamma(c))=\beta(X_0,\Gamma(0))$ for all remaining
$c\in C$.
\end{proposition}

\begin{proof}
The integral threefold $Y$ dominates the smooth curve, so it is flat:
its local rings are torsion-free modules over the corresponding
discrete valuation rings.  Its fibres have pure dimension two, and
it is projective over $C$.  Lemma~\ref{lem:square-specialization} and
the assumed vanishing give, for $c\ne0$,
\[
 (K_{Y_c}+B_c)^2
 = (L|_{Y_c})^2
 = (L|_{S_0})^2+\sum_a(L|_{S_a})^2
 = (K_{S_0}+B_{S_0})^2.
\]
The restrictions of $L$ are nef over the respective $X_c$, since
$L$ is $h$-nef.  The full-boundary hypotheses therefore allow
Lemma~\ref{lem:projective-local-square} on both the central and
noncentral birational models.  The displayed equality and
$K_{X_c}^2=K_{X_0}^2$ yield the claimed equality of local invariants.
\end{proof}

When a finite change of curve was used to construct $Y$, this equality
still concerns the original fibre singularities: such a change
repeats individual fibres and does not alter their analytic germs.

Lemma~\ref{lem:supported-square} extends the specialization identity to
nonproper families when the divisor has proper support over the curve.
Its two proofs explain how the numerical comparison can be kept local.

\subsection{Constancy in the original analytic family}
The preceding results now establish the numerical hypothesis needed by
simultaneous resolution.  The passage through algebraic families has served
only to compare the scalar invariant; we return afterwards to the original
holomorphic family.

\begin{theorem}[Constancy of the logarithmic invariant]
\label{thm:beta-constant}
Let $F(x,t)$ be a holomorphic family of isolated hypersurface singularities
in $\mathbb C^3$ with $f_t(0)=0$ and constant positive Milnor number.
If the central normal surface germ $X_0=\{f_0=0\}$ is not log canonical,
then $\beta(X_t,0)=\beta(X_0,0)$ for every sufficiently small $t$.
\end{theorem}

\begin{proof}
By Proposition~\ref{prop:scalar-curve-reduction}, it suffices to compare
a polynomial family over a smooth algebraic curve at a marked non-log-canonical
fiber with its nearby general fibers.  Log canonicity is constant by
Lemma~\ref{lem:lc-persistence}, so all germs in the relevant connected
curve are non-log-canonical.  Apply Lemma~\ref{lem:compactification} to
obtain a projective hypersurface family with exactly these singularities
and smooth complement of the section.  Its punctured-fiber Euler
characteristic and its canonical square are constant.

Construct the model of Proposition~\ref{prop:log-model}, making its finite
smooth curve base change.  Write
\[
 Y_0=S_0+\sum_a S_a,\qquad L=K_Y+H,
 \qquad S^\circ=S\setminus\operatorname{Supp}B_S.
\]
Lemma~\ref{lem:surface-boundary} and
Proposition~\ref{prop:full-exceptional-boundary} verify all the
adjunction and full-boundary requirements.  In particular,
\[
 S_0^\circ\simeq X_0\setminus\Gamma,
 \qquad Y_c\setminus H\simeq X_c\setminus\Gamma\quad(c\ne0).
\]
The first of these open surfaces is smooth.  On every other component,
any isolated interior singularity is smoothed by the family, and
Lemma~\ref{lem:interior-correction} gives a nonnegative correction
$\chi(F_x)-1$ to its nearby Euler value.

Apply Proposition~\ref{prop:nearby-euler-comparison}.  Constancy of the
punctured-fiber Euler characteristic and the two displayed isomorphisms
cancel the contribution of $S_0^\circ$, giving
\begin{equation}\label{eq:zero-euler-sum}
 0=\sum_a\chi_c(S_a^\circ)+\delta,
 \qquad
 \delta=\sum_a\sum_{x\in\operatorname{Sing}(S_a)\cap S_a^\circ}
                  (\chi(F_x)-1)\ge0.
\end{equation}
Each $S_a$ is projective and its adjoint divisor is nef, so
Lemma~\ref{lem:surface-bmy} gives
\[
 0\le(K_{S_a}+B_{S_a})^2\le3\chi_c(S_a^\circ).
\]
Every term on the right side of \eqref{eq:zero-euler-sum} is therefore
nonnegative.  Every extra logarithmic square is zero.

All hypotheses of Proposition~\ref{prop:beta-from-log-model} have now
been verified.  It gives equality of the logarithmic invariant on the
marked and nearby general fibers.  The finite base change used in the
model repeats the individual surface germs, so this is equality for
the original algebraic family as well.  Proposition~\ref{prop:scalar-curve-reduction}
then yields the assertion for the given holomorphic family.
\end{proof}

\section{From the numerical invariant to a resolved family}
\label{sec:resolution}

We now assume that the central surface is not log canonical and that
the numerical argument has established constancy of
$\beta(X_t,0)=-P_{t,\mathrm{loc}}^2$.  Okuma's theorem gives a smooth
resolved family after a finite disk cover, with reduced nodal exceptional
restrictions.  We then use the fixed-ball Euler identity to show that
none of their nodes can smooth.  This yields the relative normal-crossing
divisor needed for complete transport in the next section.

A surface resolution is \emph{semigood} when its reduced exceptional
divisor has normal crossings; irreducible components may have self-nodes.
For a divisor in a smooth family of surfaces, \emph{relative normal
crossings} means that near its points there are coordinates $(z_1,z_2,u)$
in which the family map is $u$ and the divisor is $z_1=0$ or $z_1z_2=0$.

\begin{proposition}\label{prop:resolved-family}
Let $F(x,t)=f_t(x)$ satisfy the hypotheses of the main theorem.  Assume
that $X_0=\{f_0=0\}$ is not log canonical and that $\beta(X_t,0)$ is
constant.  After shrinking the parameter disk and making a finite base
change $t=u^e$, there is a proper modification
\[
 b:N\longrightarrow X'=\{F(x,u^e)=0\}
\]
such that $N\to\Delta_u$ is smooth, $b$ is an isomorphism off the origin
section, and its reduced exceptional divisor $E$ has relative normal
crossings.  The restriction $b_u$ is the minimal semigood resolution,
and $E|_{N_u}$ is its full reduced exceptional divisor.
\end{proposition}

\begin{proof}
The zero-fiber family is a flat deformation of a normal Gorenstein
surface singularity.  Each fiber has just the one singular point in the
chosen representative.  Thus constancy of the invariant at the section
is also constancy of the sum over the singularities of the fiber, as
used in Okuma's notation.

We first check the small-deformation setup in
\cite[Sections~3--4]{Okuma2004}.  Resolve the analytic total space, and
then principalize the pulled-back ideal of the section on the smooth
resolution \cite[Theorems~2.0.1 and~2.0.3, Lemma~4.0.3]{Wlodarczyk2009}.
The composite is proper and isomorphic off the section: the total space
is smooth there and the pulled-back ideal is the unit ideal there.
Its full reduced inverse image of the section is a simple normal crossing
divisor.  The exceptional support is proper over
the disk.  The critical loci for the parameter map on the resolution
and on all closed intersections of exceptional components lie in this
proper support.  These intersections are smooth.  On each of them a
critical component has constant parameter value; an intersection
contained in a fiber also has a single value.  Properness makes their
images closed analytic subsets of the disk.  After shrinking, their
only possible value is zero.  Hence the resolution over the punctured
disk is smooth, with relative normal crossing exceptional divisor.
These relative normal-crossing charts give a locally trivial deformation
of the exceptional divisor there.  This supplies the punctured-disk weak
simultaneous resolution required by Okuma's setup.

Okuma's Theorem~1.2 \cite{Okuma2004} now gives a simultaneous resolution
after a finite base change, with minimal semigood fiber resolutions and
a reduced divisor $E$ whose every restriction is the full reduced
exceptional divisor over the distinguished non-log-canonical section.
That section is our origin section, because there are no other
singularities.  Take the normalization of a branch of the finite base
change.  A local parameter then puts its map to the original disk in
the form $t=u^e$.  Flatness and smooth fibers make the resulting
resolved family smooth over this disk.  Formation of the Cartier
divisor restrictions commutes with this further base change.  In
particular their reducedness also holds at $u=0$.  The divisor $E$
contains no fiber, and its restrictions are effective Cartier
divisors; hence it is a relative effective Cartier divisor.  Its
support is exactly $b^{-1}(\{0\}\times\Delta_u)$.

It remains to prove relative normal crossings, since fiberwise normal
crossings alone do not give that conclusion.  Choose the fixed closed
Milnor ball $B$ from the preliminary Euler calculation and set
$N_B=b^{-1}(X'\cap(B\times\Delta_u))$.  Its projection to the disk
is a proper smooth submersion with boundary.  Near the sphere boundary
the modification is an isomorphism, and the required submersion follows
from the uniform boundary transversality of the original family.
Consequently $\chi((N_B)_u)$ is constant.  The complement of $E_u$ in
this fiber is isomorphic to $(X_{u^e}\cap B)\setminus\{0\}$.  The
preliminary identity $\chi(X_t\cap B)=1$ and additivity, using compact
supports on the complements, give
\[
 \chi(E_u)=\chi((N_B)_u),
\]
so the compact reduced nodal curves $E_u$ have constant Euler
characteristic.

Here is the local consequence of that constancy.  The divisor $E$ is
Cartier in the smooth total space $N$.  At a smooth point of $E_0$ it
is already smooth over the parameter.  At each node of $E_0$, the
holomorphic Morse lemma with parameter gives relative coordinates in
which its equation is
\[
 z_1z_2=a(u),\qquad a(0)=0.
\]
There are finitely many nodes.  Choose disjoint small neighborhoods of
them with transverse boundaries.  Smoothness over the parameter near
the remaining points of $E_0$, together with properness of $E\to\Delta_u$,
excludes new nodes outside these neighborhoods after shrinking the disk.
The remainder is a proper smooth family with boundary, so its topology
is unchanged.  Inside a node neighborhood a crossing has Euler
characteristic one, whereas a smoothing is an annulus and has Euler
characteristic zero.  Thus for general small $u$,
\[
 \chi(E_u)=\chi(E_0)-
 \#\{\text{nodes whose smoothing function }a\text{ is not identically zero}\}.
\]
No global labeling of components is needed; the calculation includes
self-intersection nodes.  Since the Euler characteristic is constant,
every smoothing function vanishes identically.  The resulting charts,
together with the charts at the smooth points, prove relative normal
crossings near $E_0$.  Properness of $E\to\Delta_u$ extends this
conclusion over a smaller disk.
\end{proof}

We have obtained the relative normal-crossing divisor needed for complete
transport.  The next section constructs the ambient homeomorphisms and
then removes the finite parameter cover.

\section{From resolutions to ambient right triviality}
\label{sec:ambient-lifting}

This section supplies the final topological step. Its input is complete
transport on the punctured zero fibers; its output is a trivialization of
the functions on ambient neighborhoods. We first make that passage on fixed
closed-ball tubes. We then construct the required transport from either
Whitney regularity or a simultaneous resolution, and finally remove a finite
base change. No estimate for a vector field at the singular section is needed.

\begin{lemma}[Ambient lifting on a closed-ball tube]
\label{lem:ambient-tube}
Let $G(x,u)$ be holomorphic on a neighborhood of
$B\times D$, where $B\subset\mathbb C^n$ is a closed ball centered at $0$
and $D\subset\mathbb C$ is a disk centered at $0$.
Suppose that $G(0,u)=0$ and that $d_xG(x,u)\ne0$ for $x\ne0$.
Assume that, for some $\eta>0$, the levels
$G(\cdot,u)=c$, $|c|<\eta$, meet $\partial B$ transversely.
Put
\[
 T_u=\{x\in B:|G(x,u)|<\eta\},\qquad
 Z^*=\{(x,u)\in(B\setminus\{0\})\times D:G(x,u)=0\}.
\]
Suppose there are smooth real vector fields $v_1,v_i$ on $Z^*$,
tangent to its sphere boundary, projecting to $1,i$ on $D$.
Suppose moreover that their transport along each segment $[0,u]$
and its reverse exists for the entire segment in $Z^*$.
More precisely, this is the transport of
$\operatorname{Re}(u)v_1+\operatorname{Im}(u)v_i$ for time one,
starting over $0$, and its reverse starting over $u$.
Then, after shrinking $D$ if necessary, there are homeomorphisms
\[
 H_u:T_0\longrightarrow T_u,
 \qquad G(H_u(x),u)=G(x,0),
\]
such that $H_0=\mathrm{id}$, $H_u(0)=0$, and both the family and its
inverse are jointly continuous. Restriction to the interiors of the tubes
is a parameter-preserving ambient homeomorphism near the origin section.
\end{lemma}

\begin{proof}
On the ambient tube with the origin section removed, $(G,u)$ is a
smooth submersion, including on the sphere boundary. Extend each $v_a$,
$a\in\{1,i\}$, to a smooth ambient field $V_a$ satisfying
\[
 du(V_a)=a,\qquad dG(V_a)=0,
 \qquad V_a\text{ tangent to the sphere boundary}.
\]
Here the prescribed extension along $Z^*$ can be made in submersion
coordinates, with boundary coordinates where necessary. Equivalently,
subtract a preliminary lift and extend the remaining section of the
kernel bundle of $d(G,u)$ from the closed submanifold $Z^*$ of the axis
complement. A partition of unity patches these extensions: the specified
projections, boundary tangency, and prescribed values on $Z^*$ are affine
linear conditions. There is no assertion that $V_a$ extends to the axis.

For $u\in D$, integrate
$\operatorname{Re}(u)V_1+\operatorname{Im}(u)V_i$ for time one,
starting at $(x,0)$. For $x\ne0$ on the zero level, existence for the
whole interval is the hypothesis. For $G(x,0)=c\ne0$, the trajectory
stays in the compact set
\[
 \{(y,s):y\in B,\ s\in[0,u],\ G(y,s)=c\},
\]
which misses the axis. Smooth continuation and boundary tangency give
existence for the whole interval here as well. Reverse transport gives
the inverse on $T_u\setminus\{0\}$. Both maps depend smoothly on their
arguments off the axis, also in charts at the sphere boundary, and preserve
the value of $G$.

Set $H_u(0)=0$. To prove joint continuity there, let $x_j\to0$ and
$u_j\to u$. A subsequential limit $y$ of $H_{u_j}(x_j)$ exists in $B$,
and function preservation gives $G(y,u)=0$. If $y\ne0$, smooth dependence
of reverse transport near $(y,u)$ would imply that $x_j$ tends to its
nonzero endpoint over $0$, a contradiction. Thus every such limit is $0$.
The same argument proves inverse continuity: if $y_j\in T_{u_j}$ tends
to $0$, any nonzero subsequential limit of $H_{u_j}^{-1}(y_j)$ contradicts
smooth forward transport. Consequently the maps and inverses are jointly
continuous everywhere. Reverse transport also proves that each $H_u$ is
onto exactly $T_u$. Finally, the tube interiors form open neighborhoods of
the respective origins and their union is an open neighborhood of the
origin section. This proves the last assertion.
\end{proof}

\begin{corollary}[Whitney regular zero fibers]
\label{cor:whitney-right}
Let $G(x,u)$, with $x\in\mathbb C^n$, be a holomorphic family with
$G(0,u)=0$ and no spatial critical points off the origin section near $(0,0)$. If the pair
\[
 \bigl(\{G=0\}\setminus\{x=0\},\ \{x=0\}\bigr)
\]
satisfies Whitney conditions \textup{(a)} and \textup{(b)}, then $G$ is
ambiently topologically right trivial near that section, with the central
map equal to the identity.
\end{corollary}

\begin{proof}
Write $\rho(x,u)=\|x\|^2$. We claim that $(u,\rho)$ is a real
submersion on the punctured zero locus sufficiently near $(0,0)$.
Otherwise choose points $(x_j,u_j)\to(0,0)$ witnessing failure.
After taking a subsequence, the unit spatial vectors $x_j/\|x_j\|$
converge to $v$, the real tangent spaces of the total zero locus converge
to $L$, and the real tangent spaces of its fibers converge to $Q$.
Whitney \textup{(a)} implies that $L$ contains the two real base
directions. The fiber tangent spaces have real dimension $2n-2$ and lie
in $\ker du$, so
\[
 Q=L\cap\ker du:
\]
the inclusion holds by passage to limits, and both sides have dimension
$2n-2$. Whitney \textup{(b)}, applied to the secants from $(0,u_j)$,
gives $(v,0)\in L$, hence $(v,0)\in Q$. On the other hand, failure of
submersivity says that $d\rho$ vanishes on each fiber tangent space.
It follows that $(v,0)$ is orthogonal to $Q$, a contradiction to
$\|v\|=1$.

Take a sufficiently small closed ball and parameter disk inside this
submersion neighborhood. Smooth local right inverses to $d(u,\rho)$,
patched by a partition of unity, give lifts $v_1,v_i$ of the parameter
directions with $d\rho(v_a)=0$. Their trajectories preserve the positive
spatial radius. Along any fixed parameter segment they therefore remain
in a compact set disjoint from the axis, so transport is complete in both
directions. These fields are tangent to the sphere boundary. The zero
levels are transverse to that sphere; by compactness the same holds for
all sufficiently small nonzero levels and parameters, after shrinking.
Lemma~\ref{lem:ambient-tube} applies.
\end{proof}

\begin{proposition}[Lifts supplied by a resolution]
\label{prop:resolution-right}
Let $G(x,u)$ satisfy the analytic and transversality hypotheses of
Lemma~\ref{lem:ambient-tube}. Suppose that the zero family admits a
proper holomorphic map $b:M\to\{G=0\}$ over the parameter disk such that
\begin{enumerate}
\item $q:M\to D$ is a smooth family;
\item $b$ is an isomorphism away from the origin section;
\item $E=b^{-1}(\{x=0\})$ is a relative normal-crossing divisor.
\end{enumerate}
Here the last condition means that, in local coordinates consisting of
fiber coordinates and $u$, the support of $E$ is a union of coordinate
hyperplanes in the fiber coordinates. Then $G$ has the trivialization
of Lemma~\ref{lem:ambient-tube}.
\end{proposition}

\begin{proof}
Choose the ball and a smaller parameter disk within representatives of
$b$. Above the closed ball, $q$ is a smooth submersion with boundary:
near the sphere, $b$ is an isomorphism and the fibers meet the sphere
transversely. Relative normal-crossing coordinates give local lifts of
the two real parameter directions tangent to every branch of $E$.
Away from $E$ use submersion coordinates, and near the sphere use boundary
submersion coordinates. A partition of unity gives smooth lifts on $M$
tangent both to $E$ and to the sphere boundary.

Properness makes the inverse image of the closed ball over each compact
parameter segment compact. The smooth lifts are therefore complete along
that segment in both directions. Their flows preserve $E$ and its
complement: tangency preserves each local branch, and uniqueness applied
in both time directions prevents a trajectory in the complement from
reaching $E$. Push the fields down through the isomorphism off $E$.
These are precisely the complete fields on the punctured zero family
required by Lemma~\ref{lem:ambient-tube}.
\end{proof}

\begin{lemma}[Removal of a finite disk base change]
\label{lem:remove-base-change}
Let $F(x,t)$ be a holomorphic family with $F(0,t)=0$.
Suppose that, after the base change $t=u^e$, $e\ge1$, the family
$G(x,u)=F(x,u^e)$ has the tube trivialization $H_u$ of
Lemma~\ref{lem:ambient-tube}. Then $F$ is ambiently topologically right
trivial over the original disk, by a section-preserving homeomorphism
whose central map is the identity.
\end{lemma}

\begin{proof}
Put $\zeta=\exp(2\pi i/e)$. For real $s\ge0$ sufficiently small,
\[
 J_s=H_{\zeta s}^{-1}\circ H_s:T_0\longrightarrow T_0
\]
is a homeomorphism preserving $f_0=F(\cdot,0)$, and $J_0=\mathrm{id}$.
These compositions have exact common domains: the two target tubes
$T_s$ and $T_{\zeta s}$ coincide since $s^e=(\zeta s)^e$.
The families $J_s$ and $J_s^{-1}=H_s^{-1}\circ H_{\zeta s}$ are jointly
continuous.

For $t=r^e\exp(i\theta)$, $r>0$ and $0\le\theta\le2\pi$, define
\[
 \phi_t
 =H_{r\exp(i\theta/e)}\circ J_{r\theta/(2\pi)}.
\]
The two values at the angular seam agree, because
\[
 H_{\zeta r}\circ J_r=H_r.
\]
Figure~\ref{fig:transport-seam} displays this endpoint identity.
\begin{figure}[ht]
\centering
\begin{minipage}[c]{0.49\textwidth}
\centering
\setlength{\unitlength}{1pt}
\begin{picture}(200,120)(-5,-16)
\put(10,10){\line(1,0){75}}
\put(10,10){\line(1,2){33.5}}
\qbezier(85,10)(85,56)(43.5,77)
\qbezier[12](10,10)(42.5,28.75)(75,47.5)
\qbezier(30,10)(30,16)(27.3,20)
\put(10,10){\circle*{3}}
\put(85,10){\circle*{3}}
\put(43.5,77){\circle*{3}}
\put(75,47.5){\circle*{3}}
\put(4,3){\makebox(0,0){$0$}}
\put(92,10){\makebox(0,0)[l]{$r$}}
\put(48,82){\makebox(0,0)[l]{$\zeta r$}}
\put(85,48){\makebox(0,0)[l]{$u=re^{i\theta/e}$}}
\put(36,21){\makebox(0,0)[l]{$\theta/e$}}
\put(75,-9){\makebox(0,0){$r,\zeta r\ \longmapsto\ t=r^e$}}
\end{picture}
\end{minipage}\hfill
\begin{minipage}[c]{0.49\textwidth}
\[
\begin{CD}
 T_0 @>{J_r}>> T_0\\
 @V{H_r}VV @VV{H_{\zeta r}}V\\
 T_r @= T_{\zeta r}
\end{CD}
\]
\end{minipage}
\[
 \phi_{r^e e^{i\theta}}
   =H_{r e^{i\theta/e}}\circ J_{r\theta/(2\pi)},
 \qquad 0\leq\theta\leq2\pi.
\]
\caption{Removing the angular seam for $t=u^e$ (the sector is schematic, with $e>1$).
One turn in the $t$-disk lifts from $r$ to $\zeta r$ in the $u$-disk.
These endpoints have the same target tube.  The square records the endpoint
correction $J_r=H_{\zeta r}^{-1}H_r$; the path $s\mapsto J_s$ from
$J_0=\mathrm{id}$ interpolates that correction as $\theta$ runs from
$0$ to $2\pi$.  The two endpoint formulas for $\phi_t$ therefore agree.}
\label{fig:transport-seam}
\end{figure}

The inverse is
\[
 \phi_t^{-1}
 =J_{r\theta/(2\pi)}^{-1}\circ H_{r\exp(i\theta/e)}^{-1},
\]
and its seam identity is
$J_r^{-1}\circ H_{\zeta r}^{-1}=H_r^{-1}$.
Thus both maps glue. At $r=0$ every constituent map is the identity,
independently of $\theta$; joint continuity of the constituent families
and their inverses therefore gives joint continuity on the polar disk
with its central circle collapsed. Set $\phi_0=\mathrm{id}$.

For every $t$, the source is the same tube $T_0$ and the target is
$\{x\in B:|f_t(x)|<\eta\}$. Moreover,
\[
 f_t\circ\phi_t=f_0\circ J_{r\theta/(2\pi)}=f_0,
 \qquad \phi_t(0)=0.
\]
Restricting to the interiors gives the asserted homeomorphism germ
$(x,t)\mapsto(\phi_t(x),t)$ over the original disk.
\end{proof}

\begin{remark}
A finite map between smooth pointed disks has the form $t=u^e$ after
a coordinate change: write it as $u^e$ times a nonvanishing holomorphic
unit and absorb a holomorphic $e$th root of that unit into $u$.
Accordingly, the preceding lemma applies to a normalized disk branch of
the finite base change used for a simultaneous resolution. It requires
no deck-equivariance of that resolution or of its chosen lifts.
\end{remark}

\section{The log canonical and double-point cases}
\label{sec:low-cases}

We finish with two cases accessible through sectional Milnor numbers.
The double-point case reduces to a Whitney regular family of plane
curves after splitting off a square.  The remaining log canonical case,
of multiplicity three, is itself Whitney regular.
For a holomorphic germ in $n$ variables with
an isolated critical point, write $\mu^{(j)}$ for the Milnor number of its
restriction to a general $j$-dimensional linear subspace, and put
$\mu^*=(\mu^{(1)},\ldots,\mu^{(n)})$.  Thus $\mu^{(n)}=\mu$ and
$\mu^{(1)}=\operatorname{mult}_0(f)-1$.

We use two equisingularity theorems.  The equimultiplicity theorem of
Fern\'andez de Bobadilla and Pe\l ka states that a family of isolated
hypersurface singularities with constant finite Milnor number has constant
multiplicity \cite[Theorem~1.1]{FernandezBobadillaPelka2024}.  Teissier's
criterion states that, for a deformation of an isolated hypersurface
singularity with a section and smooth complement of the section over the
base, constancy of $\mu^*$ implies the Whitney conditions for that
complement and the section \cite[Chapter~II, Theorem~3.9]{Teissier1973}.
The absence of critical points off the section in our representatives
verifies the latter smoothness hypothesis.

\begin{proposition}\label{prop:low-cases}
Let $F(x,t)=f_t(x)$ be a holomorphic family in three spatial variables,
with $f_t(0)=0$, an isolated critical point at $0$ in every fiber, and
constant Milnor number.  If either $\operatorname{mult}_0(f_0)=2$ or the
normal surface germ $X_0=\{f_0=0\}$ is log canonical, then $F$ is
topologically right-trivial over a sufficiently small parameter disk.
\end{proposition}

\begin{proof}
Suppose first that $\operatorname{mult}_0(f_0)=2$.  Choose a spatial
direction in which the second derivative of $f_0$ is nonzero.  The implicit
function theorem solves the corresponding first-derivative equation as
$z=a(y,t)$, where $y\in\mathbb C^2$ and $a(0,t)=0$.  Translating $z$ by
$a(y,t)$ and taking a holomorphic square root of the resulting nonvanishing
quadratic coefficient gives biholomorphic coordinates, depending
holomorphically on $t$ and preserving the origin, in which
\[
 f_t=z^2+g_t(y).
\]
The Jacobian quotient is then the Jacobian quotient of $g_t$, so $g_t$ has
an isolated critical point and $\mu(g_t)=\mu(f_t)$.  Equimultiplicity for
this plane family gives constant $\mu^*(g_t)$, whose two entries are
$\operatorname{mult}_0(g_t)-1$ and $\mu(g_t)$.  Teissier's criterion and
Corollary~\ref{cor:whitney-right} give a joint
right-trivialization of $g_t$.  Extend its maps by the identity in $z$ and
conjugate by the coordinate changes at $0$ and at $t$.  This gives the
required right-trivialization of $f_t$, with central map the identity.

Now suppose $X_0$ is log canonical.  Adjunction and inversion of
adjunction identify this condition with log canonicity of
$(\mathbb C^3,X_0)$ \cite{Kawakita2007}; the different on the normal
Cartier hypersurface in a smooth ambient space is zero.  These
single-germ algebraic statements also apply to the analytic germ by its
analytic right-equivalence to a sufficiently high polynomial jet.
Put $m=\operatorname{mult}_0(f_0)$.  On the blowup of the origin the
exceptional plane has discrepancy $2-m$, hence $m\leq3$.  Since the
origin is a critical point, $m\geq2$, and the case $m=2$ is already settled.

For $m=3$, let $\widetilde X_0$ be the strict transform and $P\simeq
\mathbb P^2$ the exceptional plane.  The crepant boundary on the blowup
is $\widetilde X_0+P$.  Adjunction to $P$ shows that
$(P,C)$ is log canonical, where $C=\widetilde X_0|_P$ is the tangent
cubic with its scheme-theoretic multiplicities.  A boundary component of
a log canonical pair has coefficient at most one.  Thus $C$ is reduced.

Equimultiplicity keeps the order of $f_t$ equal to three.  Choose a
projective line transverse to $C$ at three distinct smooth points.  This
condition persists for the cubic initial forms of $f_t$ when $t$ is
small.  The corresponding plane sections, and hence general plane
sections, are ordinary triple points.  Their Milnor number is four:
the two quadratic initial forms of their partial derivatives have no
common projective zero, so their local intersection multiplicity is
$2\cdot2$.  Consequently
\[
 \mu^*(f_t)=(2,4,\mu(f_0))
\]
throughout a sufficiently small disk.  Teissier's criterion gives Whitney
regularity, and Corollary~\ref{cor:whitney-right} supplies the
right-trivialization.
\end{proof}


\subsection{Proof of Theorem~\ref{thm:main}}
If the central multiplicity is two, or the central surface is log canonical,
Proposition~\ref{prop:low-cases} proves the theorem.  Otherwise the central
surface is non-log-canonical.  Theorem~\ref{thm:beta-constant} proves
constancy of its logarithmic invariant.  Proposition~\ref{prop:resolved-family}
then supplies a smooth simultaneous resolution with relative normal-crossing
exceptional divisor after a finite disk base change.  The fixed-ball
hypotheses of Lemma~\ref{lem:ambient-tube} follow from
Lemmas~\ref{lem:local-family-properties} and~\ref{lem:fixed-ball-euler}.
Proposition~\ref{prop:resolution-right} gives the joint ambient right
trivialization after that base change, and Lemma~\ref{lem:remove-base-change}
descends it to the original disk.  The formulas preserve the parameter,
the origin section, and the defining function, with the identity at the
central parameter and jointly continuous inverses.  These are precisely
all the assertions of Theorem~\ref{thm:main}.\qed

\appendix
\section{Local constructions for the numerical comparison}
\label{app:local-numerical}

This appendix supplies the explicit marked semistable construction used
in Section~\ref{sec:log-model}, together with two variants of the
numerical comparison.  First, properly supported intersection squares
allow the specialization argument to be made on a possibly nonproper
family.  We then construct the marked model by monomial modifications.
Finally, we give an alternative algebraic-curve comparison inside a
single irreducible component of the jet locus.

\subsection{Intersection numbers with proper support}

\begin{lemma}[Specialization of a supported square]
\label{lem:supported-square}
Let $q:Y\to C$ be a dominant morphism from an integral normal
$\mathbb Q$-factorial quasi-projective complex threefold to a smooth
connected complex algebraic curve.  Let
$A=\sum_i a_iE_i$ be a rational divisor whose support is proper over
$C$.  For a fiber write its fundamental cycle as
$[Y_c]=\sum_jm_j[S_j]$.  Then the number
\[
 Q_c(A)=\sum_jm_j(A|_{S_j})^2
\]
is independent of $c$.

Here, when $S_j$ is proper, its summand is the usual projective
intersection number.  When $S_j$ is not proper, it is not a component
of $A$, and $A|_{S_j}$ is a rational Cartier divisor with proper curve
support.  Its square means the degree obtained by intersecting that
proper curve cycle with the rational line bundle $A|_{S_j}$.
\end{lemma}

\begin{proof}
The morphism $q$ is flat: its local rings are torsion-free over the
local discrete valuation rings of $C$.  Its fibers have pure dimension
two.  Every vertical $E_i$ is proper, since the support of $A$ is
proper over $C$.  A horizontal $E_i$ is projective and flat over $C$:
properness together with quasi-projectivity gives projectivity, and
flatness again follows from torsion-freeness.  All intersections below
therefore have proper support.  Multiplying the divisors by common
positive integers makes them Cartier; we divide by these integers
after taking intersection numbers.

Commutativity and additivity of Cartier-divisor intersections give
\begin{align*}
 Q_c(A)
 &=\sum_i a_i\sum_jm_j
       \deg\bigl(A\cdot E_i\cdot[S_j]\bigr)\\
 &=\sum_i a_i\deg\bigl(A|_{E_i}\cdot Y_c|_{E_i}\bigr).
\end{align*}
For distinct $E_i$ and $S_j$, the two orders of restriction describe
the same proper intersection curves with their intersection
multiplicities.  If $E_i=S_j$, all calculations take place on a proper
surface and use its normal-bundle class.  These observations also
justify the formula when some fiber component is not proper.

For a vertical $E_i$, the line bundle of $Y_c$ restricts trivially to
$E_i$, since it comes from a point divisor on $C$ and $q|_{E_i}$ is
constant.  Its term is zero, including when $E_i\subset Y_c$.
For a horizontal $E_i$, the term is the degree of the fixed rational
line bundle $A|_{E_i}$ on the scheme-theoretic fiber $(E_i)_c$.
That degree is constant in a projective flat curve family.  For
example, it is the coefficient of $n$ in the fiberwise Euler
characteristic of a Cartier multiple of $nA|_{E_i}$; flatness makes
the Euler characteristic independent of $c$, and the intersection
formula identifies this coefficient.  This proves the assertion.
\end{proof}

There is a second proof that explains why a compactification of the
ambient family causes no extra term.  Take an integral normal
projective compactification $\overline Y\to C$ of $Y\to C$.
The support of $A$ is already closed in $\overline Y$: its map to
$\overline Y$ is proper over the separated base $C$.  A Cartier
multiple of $A$ thus extends to a Cartier divisor on $\overline Y$
by gluing it to the zero divisor on
$\overline Y\setminus\operatorname{Supp}(A)$.  The compactification
is flat over $C$.  Lemma~\ref{lem:square-specialization} makes the
extended fiber square constant.  Its restriction to every component
contained in $\overline Y\setminus Y$ is trivial, and therefore has
square zero.  On a component meeting $Y$, its square is the supported
intersection number in Lemma~\ref{lem:supported-square}.  The
projective identity is consequently precisely the desired local one.

\begin{corollary}[The local exceptional-square comparison]
\label{cor:local-exceptional-square}
Suppose a projective birational morphism $h:Y\to X$ is a log model
over a smooth curve as in Section~\ref{sec:full-boundary}, except
that $X\to C$ need not be proper.  Assume $X$ and $Y$ are
quasi-projective, all exceptional prime divisors map into the section
$\Gamma\simeq C$, and put
\[
 A=K_Y+H-h^*K_X.
\]
If $Y_0=S_0+\sum_aS_a$ and the nearby fibers are irreducible, then
\[
 (A|_{Y_c})^2=(A|_{S_0})^2+\sum_a(A|_{S_a})^2.
\]
Under the full-boundary, adjunction, and nefness hypotheses of
Proposition~\ref{prop:beta-from-log-model}, the first two squares
are respectively $-\beta(X_c,\Gamma(c))$ and
$-\beta(X_0,\Gamma(0))$, while
\[
 (A|_{S_a})^2=(K_{S_a}+B_{S_a})^2.
\]
Thus vanishing of the extra squares proves constancy of the local
invariant without a projective compactification of $X$.
\end{corollary}

\begin{proof}
The divisor $A$ is exceptional and has proper support over the section,
because $h$ is projective.  Lemma~\ref{lem:supported-square} applies.
On each $S_a$, a divisor pulled back from $X$ has trivial line-bundle
class, giving the displayed adjunction formula.  On $S_0$ and on a
nearby fiber, pull back to a good resolution of the corresponding
surface germ.  As in Lemma~\ref{lem:projective-local-square}, log
canonicity makes the difference from the pulled-back logarithmic
canonical divisor an effective exceptional negative part, orthogonal
to the relatively nef part.  The pullback of $A$ is the numerical
exceptional representative of that nef part.  Its properly supported
square is therefore $-\beta$, by the definition of the local invariant.
This argument uses no global canonical square of $X_c$.
\end{proof}

\subsection{An explicit marked semistable construction}

We spell out the dimension-three toroidal construction underlying
the marked semistable reduction used in Proposition~\ref{prop:log-model};
compare \cite[Theorems~7.17 and~7.19]{KM1998}.  It preserves the
horizontal markings throughout the construction.

\begin{proposition}[Marked monomial reduction in dimension three]
\label{prop:marked-monomial-reduction}
Let $q:Z\to C$ be a morphism from a smooth quasi-projective complex
threefold to a smooth curve.  Let $0\in C$, write
$Z_0=\sum_i a_iZ_i$, and let $H$ be a reduced divisor with horizontal
components.  Suppose that $(Z_0)_{\mathrm{red}}+H$ is simple normal
crossing, and that away from $0$ the morphism is smooth with relative
simple normal crossing divisor $H$.
After shrinking $C$ about $0$, choose a parameter $\tau$ with no
other zero and an integer $b$ divisible by every $a_i$.  After the
finite base change $s^b=\tau$ and normalization, there is a projective
birational modification whose total space is smooth and whose special
fiber is reduced and simple normal crossing together with the
transformed horizontal marking.  The modification is an isomorphism
on every open set where the original morphism was smooth and the
marking relatively simple normal crossing.
\end{proposition}

\begin{proof}
Work componentwise if the normalized base change is disconnected.
At a point of $Z_0$, choose simple normal crossing coordinates in
which
\[
 \tau=u\prod_{i\in I}z_i^{a_i},
 \qquad u\text{ a unit},\qquad |I|\le3.
\]
The horizontal components present are among the other coordinate
hyperplanes.  An \'etale root of $u$ reduces the local calculation
to the monomial equation.  Each normalized local sheet is toric in
the vertical coordinates, with cone and lattice
\[
 \sigma_I=\mathbb R_{\ge0}^{I},\qquad
 N_I=\left\{(d_i)\in\mathbb Z^{I}:
                \sum_{i\in I}a_id_i\equiv0\pmod b\right\}.
\]
The order of the new base parameter is the integral height function
$\ell(d)=\sum_i a_id_i/b$.  Since $a_i\mid b$, the primitive
generator $(b/a_i)e_i$ of each original ray has height one.
The remaining coordinates, including the horizontal markings, are
unchanged product factors.

Subdivide the cones by inserting all their height-one lattice rays.
There are finitely many: their intersections with $\ell=1$ lie in
compact simplices.  At each stage insert a point not already a ray
by a star subdivision of its smallest containing cone.  Starting
with simplicial cones, this keeps every cone simplicial.  To make
the choices compatible across all charts, carry out the sequence on
the single orthant with one coordinate for each global $Z_i$ and
the congruence lattice with the same coefficients $a_i$.  Restrict
its subdivisions to each face that actually occurs.  A point not
on such a face does not change that face.

Here is a global realization by projective modifications.  A vertical
ray $v$ records the orders along its divisor of the pullbacks of the
original $Z_i$.  Let $G_v$ be the sum of all vertical prime divisors
with that order vector.  On every current chart the fan is simplicial,
so $G_v$ is rational Cartier.  Suppose a ray to be inserted is in
the relative interior of the cone generated by $v_1,\ldots,v_k$,
and write $v=\sum_j\lambda_jv_j$, with all $\lambda_j>0$ rational.
Choose positive integers $c_j$ such that the $c_jG_{v_j}$ are
Cartier and all products $c_j\lambda_j$ are equal.  Normalize the
blowup of the coherent ideal
\[
 \mathcal I=\sum_j\mathcal O(-c_jG_{v_j}).
\]
On a toric chart its generators are monomials.  The regions on which
one of their orders is smallest give exactly the star subdivision
at $v$: the equality of these orders is
$c_1d_1=\cdots=c_kd_k$ in ray coordinates.  On a chart missing one
of the required rays the corresponding summand is the unit ideal,
so the modification is the identity.  Thus these normalized blowups
realize the compatible subdivisions globally.  They are projective,
and preserve the product coordinates for the horizontal markings.

We verify regularity, including the small-dimensional lattice step.
In an actual chart the vertical cone has dimension at most three.
Its section at height one is now a segment, a triangle, or a point
with no lattice points other than its vertices.  An empty lattice
segment is primitive.  An empty lattice triangle is unimodular:
translate one vertex to zero and denote its edge vectors by $v,w$.
If their lattice index exceeded one, a nonzero lattice point of the
half-open parallelogram would have the form
$\alpha v+\gamma w$, with $0\le\alpha,\gamma<1$.  If
$\alpha+\gamma\le1$, this point lies in the triangle and is not a
vertex.  If $\alpha+\gamma>1$, the lattice point
$(1-\alpha)v+(1-\gamma)w$ does so.  Both contradict emptiness.
The edge vectors therefore form a basis of the height-zero lattice.
Adjoining one height-one vertex gives a basis of the full lattice,
so the corresponding cone is regular.  The same conclusion follows
from the primitive-segment calculation in dimension two and is
immediate in dimension one.

The resulting local toric varieties are smooth.  Every vertical ray
still has height one, so the divisor of $s$ is reduced.  Regular
toric coordinates show that this divisor together with all horizontal
coordinate divisors is simple normal crossing.  The original global
labels and order vectors distinguish local boundary branches, so
the construction does not identify distinct branches of one component.
Outside the special fiber the cover is \'etale and no subdivision
acts.  At a smooth point of the original morphism over $0$ there is
only one vertical coordinate of multiplicity one; its cone needs no
subdivision.  This proves the asserted isomorphism property as well.
\end{proof}

\subsection{Algebraic curves through prescribed jets}

The constructibility lemma in Section~\ref{sec:algebraic-reduction}
is useful in its own right.  The scalar comparison also admits the
following argument that does not use constructibility of the scalar.

\begin{lemma}[Comparison within an irreducible jet component]
\label{lem:jet-component-comparison}
Let $U_\mu$ be the locally closed algebraic jet locus of
Lemma~\ref{lem:mu-jet-locus}.  Assume the algebraic constancy assertion
in Proposition~\ref{prop:scalar-curve-reduction}.  Any two points of
the same irreducible component of $U_\mu$, one of which defines a
non-log-canonical surface germ, have the same Wahl invariant.
Consequently that proposition holds without invoking
constructibility of the Wahl invariant.
\end{lemma}

\begin{proof}
Write $Z$ for the irreducible component and let $a_0,a_1\in Z$ be
the points in question.  If they coincide there is nothing to prove.
Resolve a projective closure of $Z$, obtaining a smooth irreducible
projective variety $\widetilde Z$ with a proper surjective map to
that closure.  Choose a point above each $a_i$.  If $\dim Z=1$,
use $\widetilde Z$ itself.  If $\dim Z\ge2$, general sufficiently
ample hypersurfaces through the two chosen points cut out a smooth
irreducible projective curve through them.  Indeed the systems can
be chosen to separate first jets at those points; their general
members are smooth there, and Bertini applies away from the finite
base locus.  Iterating the irreducibility and smoothness assertions
of Bertini gives the required complete-intersection curve; compare
\cite{KleimanBertini}.

The complement of the inverse image of $Z$ is closed and misses the
chosen points.  Its intersection with the curve is therefore
finite.  Delete those finitely many points, and, if needed, restrict
the curve further to an affine open containing the two chosen
points.  The result is a smooth connected algebraic curve $C'$,
with a morphism to $U_\mu$ whose image contains $a_0$ and $a_1$.
Pulling back the universal polynomial gives regular coefficients
on $C'$ and constant Milnor number $\mu$.  Log canonicity is
constant by Lemma~\ref{lem:lc-persistence}; hence all its germs are
non-log-canonical.  Apply the assumed scalar assertion with any
point of $C'$ as the marked point.  Its value agrees with the
nearby general values.  The exceptional algebraic subset of a
curve is finite; shrinking an analytic disk about the marked
point avoids its other points.  The value is therefore constant
on that disk, including its center.  It follows that the invariant
is locally constant throughout the connected curve $C'$, and in
particular has the same value at the two prescribed points.

Finally, the finitely many irreducible components of $U_\mu$ not
containing $a_0$ form a closed subset missing $a_0$.  A sufficiently
small analytic neighborhood of $a_0$ avoids that subset.  Every
nearby jet $a(t)$ in the holomorphic jet map therefore shares an
irreducible component with $a_0$.  The comparison just proved makes
its Wahl invariant equal to the central value.  Individual finite
determinacy then returns this equality to the original holomorphic
family, exactly as in Proposition~\ref{prop:scalar-curve-reduction}.
\end{proof}

\section{Continuous gradient lifts}
\label{sec:continuous-gradient}

There is a second way to pass from transport on the punctured zero set
to ambient right triviality.  It uses vector fields that extend
continuously to the origin section, together with a condition preventing
their trajectories from reaching that section in finite time.  The
gradient lift and its correction near the zero set belong to the
vector-field method of Kuo and Parusi\'nski; see
\cite[Section~2, especially equation~(2.5)]{Parusinski1999}.
We give the criterion with all its transport hypotheses explicit.

Let $G$ be holomorphic on a neighborhood of $\{0\}\times D$ in
$\mathbb C^n\times D$, where $D$ is a complex disk centered at zero,
and suppose $G(0,u)=0$.  Set
\[
 A=\{0\}\times D,\qquad X^*=\{G=0\}\setminus A.
\]
Real tangent vectors are written in complex coordinates as $(w,a)$,
with $w\in\mathbb C^n$ and $a\in\mathbb C$.  Norms are Euclidean.
All assertions below concern sufficiently small representatives.

\begin{proposition}[A continuous-lift criterion]
\label{prop:continuous-gradient-lift}
Assume that $d_xG\ne0$ off $A$ and that, at every point of $A$,
\begin{equation}
\label{eq:gradient-small-o}
 \frac{|G_u(x,u)|}{\|d_xG(x,u)\|}\longrightarrow0
 \qquad\text{as }(x,u)\longrightarrow A,\quad x\ne0.
\end{equation}
For each $a=1,i$, suppose that $X^*$ carries a smooth real vector field
$v_a=(w_a,a)$ with the following properties:
\begin{enumerate}
\item $dG(v_a)=0$ and $w_a(x,u)\to0$ on approach to every point of $A$;
\item no trajectory of $v_a$ in $X^*$ has a finite-time endpoint in
      $A$ lying inside the representative, in either time direction.
\end{enumerate}
Then, after shrinking, there is an ambient homeomorphism germ
\[
 H(x,u)=(H_u(x),u),\qquad
 G(H_u(x),u)=G(x,0),\qquad H_u(0)=0,
\]
with $H_0=\mathrm{id}$ and with $H$ and $H^{-1}$ jointly continuous.
\end{proposition}

\begin{proof}
On the ambient complement of $A$, put
\[
 k_a=-aG_u\,
       \frac{(\overline{G_{x_1}},\ldots,\overline{G_{x_n}})}
            {\|d_xG\|^2},\qquad K_a=(k_a,a).
\]
These smooth real fields satisfy $dG(K_a)=0$.
Equation~\eqref{eq:gradient-small-o} says precisely that $k_a\to0$ at
$A$.  On $X^*$, the difference $s_a=w_a-k_a$ is a section of the
complex vector bundle $K=\ker d_xG$, and it too tends to zero at $A$.

We extend this correction while preserving its vanishing at the
section.  The smooth submanifold $X^*$ is closed in the ambient
complement of $A$.  Choose a smooth tubular retraction $r:U\to X^*$,
with a variable tube radius small enough that
\[
 \|z-r(z)\|<\tfrac12\operatorname{dist}(r(z),A)
 \qquad(z\in U).
\]
Choose a smooth function $\chi$ on the axis complement, equal to one
on $X^*$, taking values in $[0,1]$, and with support in $U$.  If
$\pi_{K_z}$ denotes Hermitian orthogonal projection onto $K_z$, define
\[
 \widetilde s_a(z)=\chi(z)\pi_{K_z}(s_a(r(z)))\quad(z\in U),
 \qquad \widetilde s_a=0\quad\text{outside }U.
\]
This is a smooth section of $K$ agreeing with $s_a$ on $X^*$.
If $z$ approaches a point of $A$ within $U$, the displayed tube bound
implies that $r(z)$ approaches the same point.  Since the projection
has norm at most one, $\widetilde s_a(z)\to0$ as well.  Thus
\[
 W_a=(k_a+\widetilde s_a,a)
\]
extends continuously across $A$ by $W_a=(0,a)$.  It is smooth off
$A$, agrees with $v_a$ on $X^*$, and annihilates $dG$ everywhere.
At points of $A$ the last assertion follows from $G(0,u)=0$.

Peano's existence theorem gives local solutions of each real ODE
$\dot z=W_a(z)$.  They are unique off $A$ because the field is
smooth there.  A solution meeting $A$ has $G=0$ throughout, by the
chain rule.  If it left $A$, a connected interval on which it lies
off $A$ would give a trajectory of $v_a$ with a finite-time endpoint
on $A$, contrary to the hypothesis.  The same argument in reverse
time excludes arrival from outside $A$.  On $A$ the equation is
$\dot x=0$, $\dot u=a$, so uniqueness holds there too.

For clarity, continuity of these local flows also holds at $A$.
On a fixed compact neighborhood the continuous field is bounded.
Initial points in a smaller neighborhood therefore have solutions on
a common short interval, all contained in that compact neighborhood.
For a convergent sequence of initial points, the solutions are
uniformly bounded and equicontinuous.  Every uniformly convergent
subsequence satisfies the integral equation for the limiting initial
point.  Uniqueness identifies its limit, and hence the whole sequence
converges.  The same argument applies to negative times.  Uniqueness
and reverse flow give continuous local inverse maps.  Consequently
these are local flows by homeomorphisms, jointly continuous in the
initial point and time.

Write $u=s+it$.  Starting at $(x,0)$, first follow $W_1$ for time
$s$ and then $W_i$ for time $t$.  Uniform short-time existence permits
both steps for all $x$ and $u$ in common smaller neighborhoods.  The
endpoint has parameter $u$ and defines $H_u(x)$.  Each step preserves
$G$ and the section.  Reverse the two steps, in the opposite order,
to obtain the joint inverse.  The reverse operations are defined on
an open neighborhood of the section.  Within that neighborhood the
image of our open source is its inverse image under the reverse map,
so this image is open.  At $u=0$ both times are zero.  This
proves all the assertions.
\end{proof}

\begin{remark}[Radial transport]
\label{rem:continuous-radial}
If the no-finite-time-endpoint condition holds for every constant real
combination of $v_1$ and $v_i$, one may instead integrate
$\operatorname{Re}(u)W_1+\operatorname{Im}(u)W_i$ for time one.
The preceding uniqueness argument applies to this combined field.
Its speed on a compact neighborhood is bounded by a constant times
$|u|$, so the solutions exist up to time one after shrinking $u$ and
the initial neighborhood.  Applying the same compactness argument to
the integral equations, now also allowing $u$ to vary, gives joint
continuity.  Reversing the trajectory gives the inverse.  Thus this
stronger trajectory hypothesis provides transport along $[0,u]$.
\end{remark}

The small-o condition and the trajectory condition have different
roles.  The L\^e--Saito--Teissier criterion \cite{LeSaito1973} supplies
\eqref{eq:gradient-small-o} for the $\mu$-constant isolated families
considered here; see \cite[Chapter~II, Remark~3.10]{Teissier1973} and
\cite[proof of Corollary~2.1, equation~(2.6)]{Parusinski1999}.
The latter writes the estimate with the full gradient; it is
equivalent to the spatial-gradient version since
$\|dG\|^2=\|d_xG\|^2+|G_u|^2$.
The following two constructions supply the additional trajectory
condition.

\begin{corollary}[An explicit Whitney lift]
\label{cor:gradient-whitney}
Suppose that $G$ satisfies the analytic and gradient hypotheses of
Proposition~\ref{prop:continuous-gradient-lift} and that $(X^*,A)$
satisfies the Whitney conditions.  Then $G$ has the trivialization
of that proposition.  The fields on $X^*$ can be chosen to preserve
$\|x\|^2$, and every constant real combination has the same property.
\end{corollary}

\begin{proof}
Use the Hermitian inner product linear in its first entry, and set
$q(x,u)=\pi_{K_{(x,u)}}x$.  We first show that
\begin{equation}
\label{eq:whitney-projection}
 \frac{\|q(x,u)\|}{\|x\|}\longrightarrow1
 \qquad\text{on }X^*\text{ approaching }A.
\end{equation}
Along any sequence, pass to a subsequence on which $x/\|x\|$ and
$K_{(x,u)}$ converge, to $v$ and $Q$ respectively.
Equation~\eqref{eq:gradient-small-o} implies that the tangent spaces
of the total zero set converge to $Q\times\mathbb C$.
Whitney condition~(b), applied to the secants between $(x,u)$ and
$(0,u)$, puts $(v,0)$ in this limiting space.  Thus $v\in Q$, which
proves \eqref{eq:whitney-projection}.

In particular $q\ne0$ sufficiently near the section.  With $k_a$
as in the preceding proof, define
\[
 w_a=k_a-
      \frac{\operatorname{Re}\langle k_a,x\rangle}{\|q\|^2}\,q,
 \qquad v_a=(w_a,a),\qquad a=1,i.
\]
The correction belongs to $K$, so $dG(v_a)=0$.
Moreover $\operatorname{Re}\langle q,x\rangle=\|q\|^2$, giving
$d\|x\|^2(v_a)=0$.  Its norm is bounded by
$\|k_a\|\|x\|/\|q\|$; hence $w_a\to0$ by
\eqref{eq:gradient-small-o} and \eqref{eq:whitney-projection}.
Every constant real combination preserves the positive radius of a
trajectory in $X^*$, so no such trajectory can reach $A$ in finite
time.  Proposition~\ref{prop:continuous-gradient-lift} and
Remark~\ref{rem:continuous-radial} apply.
\end{proof}

\begin{corollary}[Continuous lifts from a resolution]
\label{cor:gradient-resolution}
Suppose that $G$ satisfies the analytic and gradient hypotheses of
Proposition~\ref{prop:continuous-gradient-lift}.  Suppose that there
is a proper holomorphic map $b:M\to\{G=0\}$ over $D$, an
isomorphism away from $A$, such that the composite $q:M\to D$ is
smooth and $E=b^{-1}(A)$ is a relative
normal-crossing divisor.  Then $G$ has the trivialization of that
proposition.  The fields on $X^*$ satisfy the trajectory condition
for every constant real combination.
\end{corollary}

\begin{proof}
Relative normal-crossing coordinates give smooth lifts of the real
parameter directions $1,i$ tangent to every local branch of $E$.
Submersion coordinates supply lifts away from $E$.  A partition of
unity patches these to smooth fields $\widetilde v_1,\widetilde v_i$
on $M$ with the same projections and tangencies.  Through the
isomorphism off $E$, push them to fields $v_a=(w_a,a)$ on $X^*$.

These spatial parts tend to zero at $A$.  Indeed, properness supplies
convergent subsequences of the unique preimages of any sequence in
$X^*$ tending to a point of $A$.  Every limit lies in $E$.  The
spatial component of $b$ vanishes on $E$, so its differential on a
field tangent to a branch of $E$ is zero there.  Continuity of
$db(\widetilde v_a)$ now proves $w_a\to0$.

To verify the trajectory condition, suppose that an integral arc of
a constant real combination of the $v_a$ had a finite-time endpoint
in $A$ inside the representative.  Its lift to $M\setminus E$
solves the corresponding smooth combined field.  A terminal segment
of the downstairs arc has compact closure, whose inverse image is
compact by properness.  Smooth ODE continuation therefore extends
the lifted arc to a point of $E$.  Tangency preserves $E$, and
uniqueness in both time directions prevents an arc outside $E$ from
reaching it.  This contradiction proves the required condition.
The continuous-lift criterion applies.
\end{proof}

For the families in the main theorem, these criteria give an alternative
construction of the ambient maps.  In the double-point case of
Section~\ref{sec:low-cases}, apply the continuous Whitney criterion to
the plane family $g_t$ after splitting off a square, and then extend
the maps by the identity in the remaining coordinate and conjugate back
as in that section.
The log canonical case of multiplicity three uses the Whitney criterion
directly.  In the non-log-canonical case, use the resolved family of
Proposition~\ref{prop:resolved-family} after its finite disk cover.
The L\^e--Saito--Teissier estimate applies to each of these
$\mu$-constant families.  The closed-ball construction in
Section~\ref{sec:ambient-lifting} also supplies exact common tube
domains, which are used in Lemma~\ref{lem:remove-base-change} to
descend from the finite cover to the original disk.

\bibliographystyle{plain}
{\raggedright\bibliography{references}}
\end{document}
